% Generated by the blueprint tool (blueprint-pdf/1) from a blueprint snapshot; edit the
% Markdown sources, not this file.  Compile with `latexmk -pdf` or by running
% `pdflatex` twice (the second run settles the contents and the page references).
% Images, if any, are read from assets/ next to this file.
\documentclass[11pt]{article}
\usepackage[T1]{fontenc}
\usepackage[utf8]{inputenc}
\usepackage{lmodern}
\usepackage{textcomp}
\usepackage[a4paper,margin=2.5cm,headheight=14pt]{geometry}
\usepackage{amsmath,amssymb,amsthm}
\usepackage{amscd}  % KaTeX's CD diagrams; part of amsmath, so always there
\usepackage{mathtools}
\usepackage{bm}
\usepackage{graphicx}
\usepackage{xcolor}
\usepackage{array,booktabs,tabularx}
\usepackage{fancyvrb}
\usepackage[normalem]{ulem}
\usepackage{microtype}
% Optional, and only with their fonts: a small TeX installation has the
% packages without the fonts, which fails at the first formula.
\IfFileExists{mathrsfs.sty}{\IfFileExists{rsfs10.tfm}{\usepackage{mathrsfs}}{}}{}
\IfFileExists{stmaryrd.sty}{\IfFileExists{stmary10.tfm}{\usepackage{stmaryrd}}{}}{}
\IfFileExists{cancel.sty}{\usepackage{cancel}}{}
\usepackage[hyphens]{url}
\usepackage[pdfusetitle,bookmarksnumbered=false]{hyperref}
\usepackage{bookmark}
\definecolor{bplink}{HTML}{1F4E9C}
\definecolor{bpmuted}{HTML}{5F676E}
\hypersetup{colorlinks=true,linkcolor=bplink,urlcolor=bplink,citecolor=bplink}
\setlength{\parindent}{0pt}
\setlength{\parskip}{0.55em plus 0.2em minus 0.1em}
% Lists as tight as the page's.
\makeatletter
\def\@listi{\leftmargin\leftmargini\parsep 0pt\topsep 0.3em plus 0.1em\itemsep 0.1em plus 0.1em}
\let\@listI\@listi
\def\@listii{\leftmargin\leftmarginii\labelwidth\leftmarginii\advance\labelwidth-\labelsep\parsep 0pt\topsep 0.1em\itemsep 0.1em}
% An ordered list that does not start at 1.
\newcommand{\bpstartat}[1]{\setcounter{\@enumctr}{\numexpr#1-1\relax}}
\makeatother
\setcounter{tocdepth}{3}
\emergencystretch=3em
\allowdisplaybreaks

% A numbered statement, set as a paper sets one: "Theorem 2.1 (Title)."
% #1 is the lead ("Theorem~2.1"), #2 what follows it in roman.
\newenvironment{bpstatement}[2]{\par\addvspace{0.6\baselineskip}\noindent\textbf{#1}#2\textbf{.}\enspace\ignorespaces}
  {\par\addvspace{0.4\baselineskip}}
\newenvironment{bpstatementit}[2]{\begin{bpstatement}{#1}{#2}\itshape\ignorespaces}{\end{bpstatement}}
% A step: a relation with prose of its own, read under the object it starts from.
\newenvironment{bpstep}[1]{\par\addvspace{0.3\baselineskip}\begin{list}{}{\leftmargin=1.6em\rightmargin=0pt\topsep=0pt}\item\relax
  {\small\color{bpmuted}#1}\par\nopagebreak}{\end{list}\par}
% Notes about where an object is written out.
\newenvironment{bpnote}{\par\small\color{bpmuted}\itshape}{\par}
\newcommand{\bpomitted}[1]{\par{\small\color{bpmuted}\itshape #1\par}}
\newcommand{\bpbodyheading}[1]{\par\addvspace{0.4\baselineskip}\noindent\textbf{#1}\par\nopagebreak}
\newcommand{\bprule}{\par\medskip\noindent\hfil\rule{0.3\linewidth}{0.4pt}\hfil\par\medskip}
\newcommand{\bpboundary}[1]{\par\nopagebreak{\small\color{bpmuted}#1\par}\nopagebreak}
% A reference to something this export leaves out, and one that resolves nowhere.
\newcommand{\bpout}[1]{#1\textsuperscript{\hyperlink{bp-outside}{\ensuremath{\ast}}}}
\newcommand{\bpbroken}[1]{\textcolor{red!70!black}{[#1?]}}
\newcommand{\bpslug}[1]{\texttt{#1}}
% A bibliography entry: label, anchor, entry; the label hangs in the margin.
\newcommand{\bpbibitem}[3]{\par\setbox0\hbox{[#1]\enspace}\ifdim\wd0<5em\setbox0\hbox to 5em{[#1]\hfil}\fi
  \hangindent=\wd0 \hangafter=1 \noindent\box0\phantomsection\label{#2}#3\par\smallskip}
% A character the export knows no way to set.
\newcommand{\bpglyph}[1]{\fbox{\scriptsize #1}}
% Images, at their own size unless they are wider than the line.
\newsavebox{\bpimagebox}
\newcommand{\bpimage}[1]{\sbox{\bpimagebox}{\includegraphics{#1}}%
  \ifdim\wd\bpimagebox>\linewidth\resizebox{\linewidth}{!}{\usebox{\bpimagebox}}\else\usebox{\bpimagebox}\fi}
\newcommand{\bpmissingimage}[1]{\fbox{\parbox{0.8\linewidth}{\centering\small\color{bpmuted}Image not included: #1}}}
\DefineVerbatimEnvironment{bpverbatim}{Verbatim}{fontsize=\small,xleftmargin=1em}
\renewcommand{\qedsymbol}{\ensuremath{\blacksquare}}
\hypersetup{pdftitle={Weil Conjectures — 11 The Lefschetz–Verdier formula},pdfsubject={Chapter 11 The Lefschetz–Verdier formula; all levels},pdfcreator={blueprint}}
% KaTeX commands the website knows
\providecommand{\R}{\mathbb{R}}\providecommand{\N}{\mathbb{N}}\providecommand{\Z}{\mathbb{Z}}
\providecommand{\Q}{\mathbb{Q}}\providecommand{\C}{\mathbb{C}}
\providecommand{\Reals}{\mathbb{R}}\providecommand{\reals}{\mathbb{R}}\providecommand{\Complex}{\mathbb{C}}
\providecommand{\cnums}{\mathbb{C}}\providecommand{\natnums}{\mathbb{N}}\providecommand{\Bbb}[1]{\mathbb{#1}}
\providecommand{\bold}[1]{\mathbf{#1}}\providecommand{\frak}[1]{\mathfrak{#1}}\providecommand{\mathbbm}[1]{\mathbb{#1}}
\providecommand{\lang}{\langle}\providecommand{\rang}{\rangle}\providecommand{\bra}[1]{\langle #1\rvert}
\providecommand{\ket}[1]{\lvert #1\rangle}\providecommand{\braket}[1]{\langle #1\rangle}
\providecommand{\set}[1]{\{#1\}}\providecommand{\Set}[1]{\left\{#1\right\}}
\providecommand{\infin}{\infty}\providecommand{\isin}{\in}\providecommand{\plusmn}{\pm}\providecommand{\sub}{\subset}
\providecommand{\sube}{\subseteq}\providecommand{\supe}{\supseteq}\providecommand{\exist}{\exists}
\providecommand{\bull}{\bullet}\providecommand{\larr}{\leftarrow}\providecommand{\rarr}{\rightarrow}
\providecommand{\lrarr}{\leftrightarrow}\providecommand{\Larr}{\Leftarrow}\providecommand{\Rarr}{\Rightarrow}
\providecommand{\Lrarr}{\Leftrightarrow}\providecommand{\uarr}{\uparrow}\providecommand{\darr}{\downarrow}
\providecommand{\harr}{\leftrightarrow}\providecommand{\lArr}{\Leftarrow}\providecommand{\rArr}{\Rightarrow}
\providecommand{\hArr}{\Leftrightarrow}\providecommand{\alef}{\aleph}\providecommand{\alefsym}{\aleph}
\providecommand{\real}{\Re}\providecommand{\image}{\Im}\providecommand{\weierp}{\wp}
\providecommand{\Alpha}{\mathrm{A}}\providecommand{\Beta}{\mathrm{B}}\providecommand{\Epsilon}{\mathrm{E}}
\providecommand{\Zeta}{\mathrm{Z}}\providecommand{\Eta}{\mathrm{H}}\providecommand{\Iota}{\mathrm{I}}
\providecommand{\Kappa}{\mathrm{K}}\providecommand{\Mu}{\mathrm{M}}\providecommand{\Nu}{\mathrm{N}}
\providecommand{\Omicron}{\mathrm{O}}\providecommand{\Rho}{\mathrm{P}}\providecommand{\Tau}{\mathrm{T}}
\providecommand{\Chi}{\mathrm{X}}\providecommand{\omicron}{o}\providecommand{\gt}{>}\providecommand{\lt}{<}
\providecommand{\And}{\&}\providecommand{\hdashline}{\hline}\providecommand{\operatornamewithlimits}[1]{\operatorname*{#1}}
\providecommand{\llbracket}{[\![}\providecommand{\rrbracket}{]\!]}\providecommand{\cancel}[1]{#1}
\providecommand{\mathscr}{\mathcal}
% the characters beyond ASCII this document uses
\DeclareUnicodeCharacter{2160}{\textup{I}}
\DeclareUnicodeCharacter{2161}{\textup{II}}
% the project's macros ([katex.macros] in blueprint.toml)
\providecommand{\A}{}\renewcommand{\A}{\mathbf A}
\providecommand{\AEn}{}\renewcommand{\AEn}{A_{E,n}}
\providecommand{\Aaf}{}\renewcommand{\Aaf}{\mathbf A}
\providecommand{\Ab}{}\renewcommand{\Ab}{\mathsf{Ab}}
\providecommand{\AffSch}{}\renewcommand{\AffSch}{\operatorname{AffSch}}
\providecommand{\Aone}{}\renewcommand{\Aone}{\mathbf A^1}
\providecommand{\Bl}{}\renewcommand{\Bl}{\operatorname{Bl}}
\providecommand{\Br}{}\renewcommand{\Br}{\mathrm{Br}}
\providecommand{\C}{}\renewcommand{\C}{\mathbb C}
\providecommand{\Cech}{}\renewcommand{\Cech}{\check{H}}
\providecommand{\Coext}{}\renewcommand{\Coext}{\operatorname{Coext}}
\providecommand{\Cone}{}\renewcommand{\Cone}{\operatorname{Cone}}
\providecommand{\Cons}{}\renewcommand{\Cons}{\operatorname{Cons}}
\providecommand{\Cref}{}\renewcommand{\Cref}[1]{\text{#1}}
\providecommand{\D}{}\renewcommand{\D}{D}
\providecommand{\DLAFT}{}\renewcommand{\DLAFT}{\operatorname{FT}}
\providecommand{\DLAFr}{}\renewcommand{\DLAFr}{\operatorname{Fr}}
\providecommand{\DLAQ}{}\renewcommand{\DLAQ}{{\overline{\mathbf Q}_\ell}}
\providecommand{\DLRQ}{}\renewcommand{\DLRQ}{{\overline{\mathbf Q}_\ell}}
\providecommand{\Dbc}{}\renewcommand{\Dbc}{D^b_c}
\providecommand{\Dc}{}\renewcommand{\Dc}{D_{\mathrm{cons}}}
\providecommand{\Dctf}{}\renewcommand{\Dctf}{D^b_{ctf}}
\providecommand{\Dic}{}\renewcommand{\Dic}{D_{\IndCons}}
\providecommand{\Dplus}{}\renewcommand{\Dplus}{D^{+}}
\providecommand{\End}{}\renewcommand{\End}{\operatorname{End}}
\providecommand{\Et}{}\renewcommand{\Et}{\mathrm{\acute{e}t}}
\providecommand{\Ex}{}\renewcommand{\Ex}{\mathrm{Ex}}
\providecommand{\Ext}{}\renewcommand{\Ext}{\operatorname{Ext}}
\providecommand{\F}{}\renewcommand{\F}{\mathbb F}
\providecommand{\FT}{}\renewcommand{\FT}{\operatorname{FT}}
\providecommand{\Fbar}{}\renewcommand{\Fbar}{{\overline{\mathbf F}_q}}
\providecommand{\Fell}{}\renewcommand{\Fell}{\mathbf F_\ell}
\providecommand{\Fib}{}\renewcommand{\Fib}{\operatorname{Fib}}
\providecommand{\Fix}{}\renewcommand{\Fix}{\operatorname{Fix}}
\providecommand{\Fp}{}\renewcommand{\Fp}{{\mathbf F_{p}}}
\providecommand{\Fq}{}\renewcommand{\Fq}{{\mathbf F_q}}
\providecommand{\Fqn}{}\renewcommand{\Fqn}{{\mathbf F_{q^{n}}}}
\providecommand{\Frob}{}\renewcommand{\Frob}{\operatorname{Frob}}
\providecommand{\GL}{}\renewcommand{\GL}{\operatorname{GL}}
\providecommand{\Gal}{}\renewcommand{\Gal}{\operatorname{Gal}}
\providecommand{\Hens}{}\renewcommand{\Hens}{\operatorname{Hens}}
\providecommand{\Hilb}{}\renewcommand{\Hilb}{\operatorname{Hilb}}
\providecommand{\Hom}{}\renewcommand{\Hom}{\operatorname{Hom}}
\providecommand{\IndCons}{}\renewcommand{\IndCons}{\operatorname{IndCons}}
\providecommand{\Kar}{}\renewcommand{\Kar}{\operatorname{Kar}}
\providecommand{\LT}{}\renewcommand{\LT}{\operatorname{LT}}
\providecommand{\Lie}{}\renewcommand{\Lie}{\operatorname{Lie}}
\providecommand{\Loc}{}\renewcommand{\Loc}{\operatorname{Loc}}
\providecommand{\Mod}{}\renewcommand{\Mod}{\mathsf{Mod}}
\providecommand{\N}{}\renewcommand{\N}{\mathbb{N}}
\providecommand{\Needspace}{}\renewcommand{\Needspace}{\par}
\providecommand{\O}{}\renewcommand{\O}{{\mathcal{O}}}
\providecommand{\Ob}{}\renewcommand{\Ob}{\mathrm{Ob}}
\providecommand{\Oe}{}\renewcommand{\Oe}{\mathcal O_E}
\providecommand{\PP}{}\renewcommand{\PP}{\mathbf{P}}
\providecommand{\Perf}{}\renewcommand{\Perf}{\operatorname{Perf}}
\providecommand{\Pic}{}\renewcommand{\Pic}{\mathrm{Pic}}
\providecommand{\Pm}{}\renewcommand{\Pm}{\mathbf P^{m}}
\providecommand{\Pone}{}\renewcommand{\Pone}{\mathbf P^{1}}
\providecommand{\Pqn}{}\renewcommand{\Pqn}{\mathbf P^{n}}
\providecommand{\ProFin}{}\renewcommand{\ProFin}{\operatorname{ProFin}}
\providecommand{\Proet}{}\renewcommand{\Proet}{\mathrm{pro\acute{e}t}}
\providecommand{\Q}{}\renewcommand{\Q}{\mathbb Q}
\providecommand{\Qb}{}\renewcommand{\Qb}{\overline{\mathbf Q}_{\ell}}
\providecommand{\Qbar}{}\renewcommand{\Qbar}{\overline{\mathbf Q}}
\providecommand{\QbarX}{}\renewcommand{\QbarX}{\overline{\mathbf Q}_{\ell,X}}
\providecommand{\Qbarl}{}\renewcommand{\Qbarl}{{\overline{\mathbf Q}_{\ell}}}
\providecommand{\Qell}{}\renewcommand{\Qell}{{\mathbf{Q}_{\ell}}}
\providecommand{\Ql}{}\renewcommand{\Ql}{\mathbf Q_{\ell}}
\providecommand{\Qlb}{}\renewcommand{\Qlb}{{\overline{\mathbf Q}_\ell}}
\providecommand{\Qlbar}{}\renewcommand{\Qlbar}{\overline{\mathbf Q}_\ell}
\providecommand{\R}{}\renewcommand{\R}{\mathbb{R}}
\providecommand{\RG}{}\renewcommand{\RG}{\operatorname{\mathrm{R}\Gamma}}
\providecommand{\RGamma}{}\renewcommand{\RGamma}{R\Gamma}
\providecommand{\RGammac}{}\renewcommand{\RGammac}{R\Gamma_c}
\providecommand{\RHom}{}\renewcommand{\RHom}{R\mathcal{H}om}
\providecommand{\RHomX}{}\renewcommand{\RHomX}{R\mathcal H om_X}
\providecommand{\Rat}{}\renewcommand{\Rat}{\operatorname{IsRat}}
\providecommand{\Res}{}\renewcommand{\Res}{\operatorname{Res}}
\providecommand{\Rlim}{}\renewcommand{\Rlim}{R\!\varprojlim}
\providecommand{\SIm}{}\renewcommand{\SIm}{\operatorname{SIm}}
\providecommand{\Sch}{}\renewcommand{\Sch}{\mathsf{Sch}}
\providecommand{\Sh}{}\renewcommand{\Sh}{\operatorname{Sh}}
\providecommand{\Span}{}\renewcommand{\Span}{\operatorname{span}}
\providecommand{\Spec}{}\renewcommand{\Spec}{\operatorname{Spec}}
\providecommand{\Tr}{}\renewcommand{\Tr}{\operatorname{Tr}}
\providecommand{\Trsm}{}\renewcommand{\Trsm}{\operatorname{Tr}^{\mathrm{sm}}}
\providecommand{\Wt}{}\renewcommand{\Wt}{\operatorname{Wt}}
\providecommand{\Z}{}\renewcommand{\Z}{\mathbb{Z}}
\providecommand{\Zar}{}\renewcommand{\Zar}{\mathrm{zar}}
\providecommand{\Zell}{}\renewcommand{\Zell}{{\mathbf{Z}_{\ell}}}
\providecommand{\Zl}{}\renewcommand{\Zl}{{\mathbf Z_\ell}}
\providecommand{\affet}{}\renewcommand{\affet}[1]{{#1}^{\mathrm{aff}}_{\mathrm{ét}}}
\providecommand{\affproet}{}\renewcommand{\affproet}[1]{{#1}^{\mathrm{aff}}_{\mathrm{proét}}}
\providecommand{\boxt}{}\renewcommand{\boxt}{\boxtimes_A^{\mathrm L}}
\providecommand{\cF}{}\renewcommand{\cF}{\mathcal{F}}
\providecommand{\cHom}{}\renewcommand{\cHom}{R\mathcal{H}\!om}
\providecommand{\cI}{}\renewcommand{\cI}{\mathcal{I}}
\providecommand{\cL}{}\renewcommand{\cL}{\mathcal L}
\providecommand{\cM}{}\renewcommand{\cM}{\mathcal M}
\providecommand{\cN}{}\renewcommand{\cN}{\mathcal N}
\providecommand{\cO}{}\renewcommand{\cO}{\mathcal O}
\providecommand{\cT}{}\renewcommand{\cT}{\mathcal T}
\providecommand{\checked}{}\renewcommand{\checked}{\color{red}\aftergroup\restorecolor}
\providecommand{\coeff}{}\renewcommand{\coeff}{\operatorname{coeff}}
\providecommand{\coker}{}\renewcommand{\coker}{\operatorname{coker}}
\providecommand{\colim}{}\renewcommand{\colim}{\operatorname{colim}}
\providecommand{\constructedby}{}\renewcommand{\constructedby}[1]{\par\smallskip\noindent\textcolor{purple!60!black}{\textsf{Constructed by:} }
  \footnotesize\sloppy\nolinkurl{#1}\par}
\providecommand{\deepresult}{}\renewcommand{\deepresult}[1]{\par\smallskip\noindent\textcolor{orange!45!black}{
  \textsf{Blueprint role: named deep theorem node; see the cluster proof architecture.}}
  \par\noindent\textcolor{orange!45!black}{\textsf{Mathematical source:} }#1\par}
\providecommand{\dt}{}\renewcommand{\dt}{\mathbin{\otimes^{\mathbf L}}}
\providecommand{\dw}{}\renewcommand{\dw}{\operatorname{dw}}
\providecommand{\et}{}\renewcommand{\et}[1]{{#1}_{\mathrm{ét}}}
\providecommand{\ev}{}\renewcommand{\ev}{\operatorname{ev}}
\providecommand{\fib}{}\renewcommand{\fib}{\operatorname{fib}}
\providecommand{\formalizationnote}{}\renewcommand{\formalizationnote}[1]{\begin{quote}\small\raggedright\textbf{Formalization note.} #1\end{quote}}
\providecommand{\fpqc}{}\renewcommand{\fpqc}{X_{\mathrm{fpqc}}}
\providecommand{\gammacont}{}\renewcommand{\gammacont}{\Gamma_{\mathrm{cont}}}
\providecommand{\gm}{}\renewcommand{\gm}{\mathbf G_m}
\providecommand{\hocolim}{}\renewcommand{\hocolim}{\operatorname{hocolim}}
\providecommand{\holim}{}\renewcommand{\holim}{\operatorname{holim}}
\providecommand{\hrB}{}\renewcommand{\hrB}{\@ifnextchar]{\hrEnd}{\hrStep}}
\providecommand{\hrEnd}{}\renewcommand{\hrEnd}{\@firstoftwo{}}
\providecommand{\hrStep}{}\renewcommand{\hrStep}{\@firstoftwo{\hrB}}
\providecommand{\id}{}\renewcommand{\id}{\operatorname{id}}
\providecommand{\kunFl}{}\renewcommand{\kunFl}{\mathbf F_\ell}
\providecommand{\kunZn}{}\renewcommand{\kunZn}{\mathbf Z/n\mathbf Z}
\providecommand{\kuncd}{}\renewcommand{\kuncd}{\operatorname{cd}}
\providecommand{\kundim}{}\renewcommand{\kundim}{\dim}
\providecommand{\kunotimes}{}\renewcommand{\kunotimes}[1]{\otimes^{\mathbf L}_{#1}}
\providecommand{\kunsupp}{}\renewcommand{\kunsupp}{\operatorname{Supp}}
\providecommand{\lotimes}{}\renewcommand{\lotimes}{\otimes^{\mathrm L}}
\providecommand{\m}{}\renewcommand{\m}{{\mathfrak{m}}}
\providecommand{\n}{}\renewcommand{\n}{{\mathfrak{n}}}
\providecommand{\nagataGap}{}\renewcommand{\nagataGap}[1]{\par\noindent\textbf{Formalization gap.} #1\par}
\providecommand{\openobligation}{}\renewcommand{\openobligation}[1]{\par\smallskip\noindent\textcolor{red!65!black}{\textsf{Open obligation:} }
  \footnotesize\sloppy\nolinkurl{#1}\par}
\providecommand{\p}{}\renewcommand{\p}{{\mathfrak{p}}}
\providecommand{\pr}{}\renewcommand{\pr}{\operatorname{pr}}
\providecommand{\preshv}{}\renewcommand{\preshv}[1]{\mathrm{PreShv}(#1)}
\providecommand{\pro}{}\renewcommand{\pro}[1]{\mathrm{Pro}(#1)}
\providecommand{\proet}{}\renewcommand{\proet}{X_{\mathrm{proét}}}
\providecommand{\q}{}\renewcommand{\q}{{\mathfrak{q}}}
\providecommand{\rad}{}\renewcommand{\rad}{\operatorname{rad}}
\providecommand{\rank}{}\renewcommand{\rank}{\operatorname{rank}}
\providecommand{\red}{}\renewcommand{\red}{\color{red}}
\providecommand{\restorecolor}{}\renewcommand{\restorecolor}{\color{black}}
\providecommand{\retiredobligation}{}\renewcommand{\retiredobligation}[1]{\par\smallskip\noindent\textcolor{red!75!black}{\textsf{Retired invalid route:} }
  \footnotesize\sloppy\nolinkurl{#1}\par}
\providecommand{\rk}{}\renewcommand{\rk}{\operatorname{rk}}
\providecommand{\sep}{}\renewcommand{\sep}{\textnormal{sep}}
\providecommand{\shv}{}\renewcommand{\shv}[1]{\mathrm{Shv}(#1)}
\providecommand{\spec}{}\renewcommand{\spec}[1]{\mathrm{Spec}(#1)}
\providecommand{\stacks}{}\renewcommand{\stacks}[1]{\href{https://stacks.math.columbia.edu/tag/#1}{\textnormal{Stacks, Tag~#1}}}
\providecommand{\stacksproject}{}\renewcommand{\stacksproject}[1]{{\cite[\href{https://stacks.math.columbia.edu/tag/#1}{Tag #1}]{stacks-project}}}
\providecommand{\statusnote}{}\renewcommand{\statusnote}[1]{\begin{quote}\small\raggedright\textbf{Status.} #1\end{quote}}
\providecommand{\tildeideal}{}\renewcommand{\tildeideal}[2]{\widetilde{(#1, #2)}}
\providecommand{\tildering}{}\renewcommand{\tildering}[3]{{#1}_{\widetilde{(#2, #3)}}}
\providecommand{\tr}{}\renewcommand{\tr}{\operatorname{tr}}
\providecommand{\wt}{}\renewcommand{\wt}{\operatorname{wt}}

\begin{document}
\pagestyle{plain}
\begin{center}
{\LARGE\bfseries Weil Conjectures\par}
\vspace{0.7em}
{\large Chapter 11: The Lefschetz–Verdier formula\par}
\vspace{0.5em}
{\small\color{bpmuted}Snapshot 2026-09-26, revision 8e4cb88e2a\par}
\end{center}
\vspace{0.5em}
{\small\tableofcontents}
\bigskip\hrule\bigskip

% object: chap-The-Lefschetz-Verdier-formula (blueprint/chap-The-Lefschetz-Verdier-formula/_section.md)
\section*{\texorpdfstring{11\quad The Lefschetz–Verdier formula}{11 The Lefschetz–Verdier formula}}
\phantomsection\label{bp1}%
\addcontentsline{toc}{section}{\texorpdfstring{11\quad The Lefschetz–Verdier formula}{11 The Lefschetz–Verdier formula}}
% text: prose-11-01
\phantomsection\label{bp2}%
The trace formula has to express the trace of Frobenius on \(R\Gamma_c(X,K)\) as a sum over the fixed points of Frobenius, that is, over \(X_0(\mathbf F_{q^m})\). We do not prove it for Frobenius directly. Instead we first prove a fixed point formula for every correspondence on a proper scheme. In this generality the formula is compatible with proper pushforward, and pushing forward to a point turns it into a tautology.

For a general correspondence, however, the fixed locus need not be discrete (Section \hyperref[bp26]{11.4}). The local side is then a sum of contributions, one for each connected component of the fixed locus, each computed in a neighbourhood of that component.

We use Grothendieck's formalism of cohomological correspondences, written up by Illusie in [\hyperref[bp57]{SGA5}, III]. In this formalism, a correspondence on \(X\) is a map \(c=(c_1,c_2):C\to X\times X\), and a cohomological correspondence on \(K\) is a morphism \(u:c_2^*K\to c_1^!K\). For example, for the graph of an endomorphism \(\varphi\) this is a map \(\varphi^*K\to K\). In particular, for Frobenius it is the Weil structure \(\phi_K:\operatorname{Fr}_X^*K\simeq K\) of Section \bpout{10.5}, up to the comparison of \(\operatorname{Fr}_X\) with the \(k\)-linear Frobenius in Chapter \bpout{12}. When \(c_2\) is proper, \(u\) induces an endomorphism \(T_c(u)\) of \(R\Gamma_c(X,K)\). In the same way \(\phi_K\) induces the Frobenius operator \(\operatorname{Frob}_{X,K}\) on \(R\Gamma_c(X,K)\) and on every stalk at a fixed point.

The local contributions come from duality, as \(\operatorname{Hom}(V,V)=V\otimes
V^\vee\) turns an endomorphism into a tensor whose evaluation is the trace. In this way \(u\) gives the trace class \(\operatorname{Tr}_c(u)\in H^0(\Fix(c),\omega_{\Fix(c)})\), supported on the fixed locus. Its integral over a proper open and closed part \(B\) of \(\Fix(c)\) is the local term \(\operatorname{LT}_B(u)\). The Lefschetz–Verdier formula (Theorem \hyperref[bp37]{11.13}) says that if \(X\) and \(C\) are proper over \(k\), then

\[
\operatorname{tr}\bigl(T_c(u)\mid R\Gamma(X,K)\bigr)
=\sum_B\operatorname{LT}_B(u).
\]

The proof involves no computation, because trace classes commute with proper pushforward of correspondences. After pushing forward to the identity correspondence of the point, the trace class becomes the trace of \(T_c(u)\) on the perfect complex \(R\Gamma(X,K)\) (Lemma \hyperref[bp30]{11.10}). The price of this generality, however, is that the formula does not compute the individual local terms; for Frobenius that is done in Chapter \bpout{12}.

With coefficients in a finite ring \(A\) both sides are traces of endomorphisms of perfect complexes (Definition \hyperref[bp5]{11.1}). They are not alternating sums of traces on cohomology modules, since these need not be projective. In fact, the trace is not even determined by the maps induced on cohomology (Example \hyperref[bp12]{11.4}). Moreover, finite rings such as \(\mathcal O_E/\pi_E^n\) cannot be avoided, because the trace formula over \(\mathcal O_E\) is obtained from its reductions.

Throughout, \(k_0=\mathbf F_q\), \(k=\overline{\mathbf F}_q\) and \(s=\operatorname{Spec}k\). For a scheme \(X\) over \(k\), \(a_X:X\to s\) is the structure morphism, and \(X=X_0\times_{k_0}k\) whenever a model \(X_0\) over \(k_0\) is given. From Section \hyperref[bp3]{11.1} until Section \bpout{12.5} the coefficient ring is an arbitrary finite commutative \(\ell\)-primary ring \(A\), with \(\ell\ne\operatorname{char}k\), and we write \(\omega_X=a_X^!A_s\) and \(\mathbb D_XK=\cHom_A(K,\omega_X)\) for the dualizing complex and the Verdier dual of Definition \bpout{10.2} with coefficients \(A\). Some proofs compute on small étale sites and transport the result by means of the comparison theorems for complexes of finite Tor amplitude.

% object: sec-Traces-of-perfect-complexes (blueprint/chap-The-Lefschetz-Verdier-formula/sec-Traces-of-perfect-complexes/_section.md), content hidden
\subsection*{\texorpdfstring{11.1\quad Traces of perfect complexes}{11.1 Traces of perfect complexes}}
\phantomsection\label{bp3}%
\addcontentsline{toc}{subsection}{\texorpdfstring{11.1\quad Traces of perfect complexes}{11.1 Traces of perfect complexes}}
% text: prose-11-02
\phantomsection\label{bp4}%
Both sides of the trace formula are traces of endomorphisms of perfect \(A\)-complexes, namely of \(R\Gamma_c(X,K)\) on the global side and of the stalks of \(K\) on the local side.

Over a field the trace of an endomorphism of a complex is the alternating sum of its traces on cohomology. Over \(A=\mathbf Z/\ell^2\), however, the cohomology modules need not be free, so they have no trace, and Example \hyperref[bp12]{11.4} below shows that the maps on cohomology do not determine the trace. We therefore define the trace on the complex. On a bounded complex of finite projective modules it is the alternating sum of the traces of the components, but the definition below expresses this sum through evaluation and coevaluation, so that it depends only on the morphism in \(D(A)\). In this form it is compatible with every change of rings (Proposition \hyperref[bp7]{11.2}), and this is how we pass from the trace formula over \(\mathcal O_E/\pi_E^n\) to the trace formula over \(\mathcal O_E\) and \(E\) (Section \bpout{12.5}).

% object: definition-categorical-trace (blueprint/chap-The-Lefschetz-Verdier-formula/sec-Traces-of-perfect-complexes/definition-categorical-trace/_section.md)
\phantomsection\label{bp5}%
\begin{bpstatement}{Definition~11.1}{ (Categorical trace)}
For \(P\in\operatorname{Perf}(A)\) put \(P^\vee=R\!\operatorname{Hom}_A(P,A)\). If \(v:P\to P\) is a morphism in \(D(A)\), its trace is the scalar

\[
\operatorname{tr}_A(v\mid P)=
\left[A\xrightarrow{\operatorname{coev}}P\otimes^{\mathbf L} P^\vee
\xrightarrow{v\otimes^{\mathbf L}1}P\otimes^{\mathbf L} P^\vee
\xrightarrow{\sim}P^\vee\otimes^{\mathbf L} P
\xrightarrow{\operatorname{ev}}A\right]\in A.
\]

Here \(\operatorname{coev}\) and \(\operatorname{ev}\) are the coevaluation and evaluation of the dual pair \((P,P^\vee)\) in \(D(A)\), and the unlabelled arrow is the symmetry isomorphism, which includes the Koszul sign.
\end{bpstatement}

% text: prose-11-03
\phantomsection\label{bp6}%
For \(P=A\) in degree zero and \(v\) multiplication by \(a\), coevaluation is \(1\mapsto1\otimes1\) and evaluation is multiplication, so the trace is \(a\). The next proposition lists the properties of the trace that we use. Compatibility with change of rings lets us move the trace formula between coefficient rings, and the chain-level formula is how we compute traces in examples.

% object: proposition-properties-trace (blueprint/chap-The-Lefschetz-Verdier-formula/sec-Traces-of-perfect-complexes/proposition-properties-trace/_section.md)
\phantomsection\label{bp7}%
\begin{bpstatementit}{Proposition~11.2}{ (Properties of the trace)}
The trace is \(A\)-linear in \(v\), invariant under conjugacy, and compatible with derived scalar extension: for every ring map \(A\to B\),

\[
\operatorname{tr}_A(v\mid P)\longmapsto
\operatorname{tr}_B(v\otimes^{\mathbf L} B\mid P\otimes^{\mathbf L} B).
\]

For a bounded complex of finite projective modules and a chain endomorphism \(v^\bullet\) representing \(v\) one has

\[
\operatorname{tr}_A(v\mid P)
=\sum_i(-1)^i\operatorname{Tr}_A(v^i\mid P^i).
\]

If \(P\) is the underlying complex of a two-step filtered object \(F^1P\to P\to P/F^1P\) of the filtered derived category, with \(F^1P\) and \(P/F^1P\) perfect, and \(v\) is the underlying map of an endomorphism \(\tilde v\) of this filtered object, with graded pieces \(v_1\) on \(F^1P\) and \(v_0\) on \(P/F^1P\), then

\[
\operatorname{tr}_A(v\mid P)=\operatorname{tr}_A(v_1\mid F^1P)
+\operatorname{tr}_A(v_0\mid P/F^1P).
\]

Over a field \(L\) it is also

\[
\operatorname{tr}_L(v\mid P)
=\sum_i(-1)^i\operatorname{Tr}_L(H^i(v)\mid H^i(P)).
\]
\end{bpstatementit}

% proof: proof/proposition-properties-trace
\begin{proof}[{Proof}]\phantomsection\label{bp8}%
Everything follows from the chain formula, which we prove first by writing evaluation and coevaluation in a finite dual basis. If \(e_1,\dots,e_n\) and \(e^1,\dots,e^n\) are dual bases of \(A^n\) and its dual, then \(\operatorname{coev}(1)=\sum_je_j\otimes e^j\), and the composite of the definition sends \(1\) to \(\sum_je^j(v(e_j))\), the ordinary trace of a matrix. The same computation applies to a finite projective module, which is a direct summand of such an \(A^n\). On a complex, the symmetry interchanging a term of degree \(i\) and its dual of degree \(-i\) contributes \((-1)^i\), which proves the chain formula. Its right-hand side is invariant under chain homotopy, since the supertrace of \(dh+hd\) is zero by cyclicity and a shift of index.

Tensoring a dual basis with \(B\) preserves evaluation and coevaluation. Hence the trace is compatible with extension of scalars for every ring map \(A\to B\). Linearity and invariance under conjugacy follow from the same description. For additivity, note that a two-step filtered object of the filtered derived category may be represented by a filtered complex of projective modules whose filtration is split in each degree. The endomorphism is then block triangular, so its supertrace is the sum of the supertraces on the two graded pieces. Finally, over a field the boundaries \(B^i\) and cycles \(Z^i\) are \(v\)-invariant subspaces \(B^i\subset Z^i\subset P^i\), with \(Z^i/B^i=H^i(P)\), and the differential induces a \(v\)-equivariant isomorphism \(P^i/Z^i\simeq B^{i+1}\). Hence

\[
\operatorname{Tr}(v^i\mid P^i)=\operatorname{Tr}(v\mid B^i)
+\operatorname{Tr}(H^i(v)\mid H^i(P))+\operatorname{Tr}(v\mid B^{i+1}),
\]

and in the alternating sum the terms for the \(B^i\) cancel in pairs, leaving the cohomological formula.
\end{proof}

% text: prose-11-04
\phantomsection\label{bp9}%
In the additivity statement of the proposition we ask for an endomorphism of the filtered object, and not merely of the three vertices of a triangle. The reason is that the middle component is not determined by the outer ones, and even for a morphism of triangles the trace of the middle component need not be the sum of the outer traces.

% object: remark-coherent-additivity (blueprint/chap-The-Lefschetz-Verdier-formula/sec-Traces-of-perfect-complexes/remark-coherent-additivity/_section.md)
\phantomsection\label{bp10}%
\begin{bpstatement}{Remark~11.3}{ (Additivity needs a filtered endomorphism)}
The coherence condition is needed over a general ring: three endomorphisms of a distinguished triangle whose squares commute only in the triangulated category need not satisfy trace additivity [Stacks, \href{https://stacks.math.columbia.edu/tag/087J}{087J}]. Every localization sequence we use below is obtained functorially from the Frobenius correspondence and the filtration by an open subscheme and its closed complement, and hence it lifts to the filtered derived category.
\end{bpstatement}

% text: prose-11-05
\phantomsection\label{bp11}%
The following example shows why the trace has to be taken on the complex.

% object: example-trace-not-determined-cohomology (blueprint/chap-The-Lefschetz-Verdier-formula/sec-Traces-of-perfect-complexes/example-trace-not-determined-cohomology/_section.md)
\phantomsection\label{bp12}%
\begin{bpstatement}{Example~11.4}{ (A trace not determined by cohomology)}
Let \(A=\mathbf Z/\ell^2\) and let \(P=[A\xrightarrow{\ell}A]\), with the terms in degrees \(0\) and \(1\). This is the perfect complex from Section \bpout{3.2} whose cohomology modules \(H^0(P)=\ell A\) and \(H^1(P)=A/\ell A\) are not perfect. The chain maps \(1_P=(1,1)\) and \(v=(1+\ell,1)\) commute with the differential, since \(\ell(1+\ell)=\ell\) in \(A\), and both induce the identity on \(H^0(P)\) and on \(H^1(P)\). By the chain formula of Proposition \hyperref[bp7]{11.2},

\[
\operatorname{tr}_A(1_P\mid P)=1-1=0,
\qquad
\operatorname{tr}_A(v\mid P)=(1+\ell)-1=\ell\neq0.
\]

Thus, over a finite ring, the trace of an endomorphism of a perfect complex is not determined by the maps it induces on cohomology. The trace formula has to be stated for the operator on the complex.
\end{bpstatement}


% object: sec-Correspondences (blueprint/chap-The-Lefschetz-Verdier-formula/sec-Correspondences/_section.md), content hidden
\subsection*{\texorpdfstring{11.2\quad Correspondences}{11.2 Correspondences}}
\phantomsection\label{bp13}%
\addcontentsline{toc}{subsection}{\texorpdfstring{11.2\quad Correspondences}{11.2 Correspondences}}
% text: prose-11-06
\phantomsection\label{bp14}%
As in topology, a pair of maps \(c_1,c_2:C\to X\) should act on \(R\Gamma_c(X,K)\) by pulling back along \(c_2\) and pushing forward along \(c_1\). The pullback \(R\Gamma_c(X,K)\to R\Gamma_c(C,c_2^*K)\) exists when \(c_2\) is proper. To push forward we then need a map \(Rc_{1!}L\to K\), where \(L\) is the complex on \(C\) we have arrived at. By the adjunction \(Rc_{1!}\dashv c_1^!\) there is a canonical such map exactly when \(L=c_1^!K\), namely the counit. What connects the two steps is therefore a map \(u:c_2^*K\to c_1^!K\). The resulting endomorphism \(T_c(u)\) is defined in Definition \hyperref[bp33]{11.11}.

With \(c_1^*\) in place of \(c_1^!\) the second step fails, since there may be no nonzero map \(Rc_{1!}c_1^*K\to K\). For example, for the inclusion \(i\) of a rational point \(y\) into a smooth curve \(X\),

\[
\operatorname{Hom}(i_*i^*\Lambda_X,\Lambda_X)
=\operatorname{Hom}(\Lambda_y,i^!\Lambda_X)
=\operatorname{Hom}(\Lambda_y,\Lambda_y(-1)[-2])=0,
\]

so the only map \(Ri_!i^*\Lambda_X\to\Lambda_X\) is zero. When \(c_1\) is finite étale, smooth purity gives \(c_1^!=c_1^*\), and the counit is summation over the fibres. Then \(T_c(u)\) is pullback along \(c_2\) followed by a sum over the fibres of \(c_1\), as for a Hecke operator.

The Lefschetz–Verdier formula (Theorem \hyperref[bp37]{11.13}) is stated in this generality. For the trace formula we apply it to the graph of Frobenius (Definition \bpout{12.2}).

% object: definition-correspondence (blueprint/chap-The-Lefschetz-Verdier-formula/sec-Correspondences/definition-correspondence/_section.md)
\phantomsection\label{bp15}%
\begin{bpstatement}{Definition~11.5}{ (Correspondences)}
Let \(X\) be a separated \(k\)-scheme of finite type. A \emph{correspondence} on \(X\) is a morphism of \(k\)-schemes \(c=(c_1,c_2):C\to X\times_sX\) from a separated \(k\)-scheme of finite type (similarly over a base \(S\)). For \(K\in D^b_{\mathrm{ctf}}(X,A)\), a \emph{cohomological correspondence} on \((X,K)\) is a morphism

\[
u:c_2^*K\longrightarrow c_1^!K.
\]

The \emph{fixed locus} of \(c\) is the closed subscheme

\[
\Fix(c)=C\times_{X\times X,\Delta_X}X,
\qquad i_c:\Fix(c)\longrightarrow C,\quad d_c:\Fix(c)\longrightarrow X,
\]

where \(i_c\) and \(d_c\) are the two projections.
\end{bpstatement}

% text: prose-11-07
\phantomsection\label{bp16}%
Here \(c_2\) is the leg along which \(K\) is pulled back and \(c_1\) is the leg along which we take \(c_1^!\); in [\hyperref[bp57]{SGA5}, III, §3] the two names are interchanged. Since \(X\) is separated, \(\Delta_X\) and \(i_c\) are closed immersions.

The example to keep in mind is the graph of an endomorphism \(\varphi:X\to X\), that is, \(C=X\) with \(c_1=\operatorname{id}_X\) and \(c_2=\varphi\). Then \(c_1^!K=K\), and a cohomological correspondence is a morphism \(u:\varphi^*K\to K\). In this case the fixed locus \(\Fix(c)\) is the fixed point scheme of \(\varphi\), which we also denote \(\Fix(\varphi)\). For the Frobenius \(\Phi_X\) and a Weil complex this is the correspondence whose trace formula we want.


% object: sec-The-exceptional-Kunneth-formula (blueprint/chap-The-Lefschetz-Verdier-formula/sec-The-exceptional-Kunneth-formula/_section.md), content hidden
\subsection*{\texorpdfstring{11.3\quad The exceptional Künneth formula}{11.3 The exceptional Künneth formula}}
\phantomsection\label{bp17}%
\addcontentsline{toc}{subsection}{\texorpdfstring{11.3\quad The exceptional Künneth formula}{11.3 The exceptional Künneth formula}}
% text: prose-11-08
\phantomsection\label{bp18}%
A cohomological correspondence \(u:c_2^*K\to c_1^!K\) is a morphism on \(C\), so to obtain from it a class on the fixed locus we must first view it as a section of \(c^!\) of a single complex on \(X\times X\), which can then be restricted to the diagonal. In linear algebra this is the identification \(\operatorname{Hom}(V,W)=W\otimes V^\vee\), under which the trace becomes evaluation. Its sheaf version is the isomorphism \(\cHom(p_2^*K,p_1^!K)\simeq K\boxtimes\mathbb D_XK\) of Lemma \hyperref[bp22]{11.7}. Here \(p_1,p_2:X\times X\to X\) are the projections and \(M\boxtimes N=p_1^*M\otimes^{\mathbf L} p_2^*N\) is the external tensor product in \(D((X\times X)_{\mathrm{pro}\acute{e}t},A)\). Applying \(c^!\) to it, we obtain \(\cHom(c_2^*K,c_1^!K)\). On the point the isomorphism is exactly \(\operatorname{Hom}(V,V)=V\otimes V^\vee\).

The isomorphism comes from Künneth formulas for \(Rf_!\), \(Rf_*\), \(f^!\), \(\cHom\) and Verdier duals. For \(Rf_!\) the formula follows directly from base change for compact supports and the projection formula, and for ordinary direct images from smooth and proper base change (Proposition \bpout{8.18}, Section \bpout{8.9}). The formula for \(f^!\) (Proposition \hyperref[bp19]{11.6}), by contrast, needs a dévissage on the coefficients, starting from constant coefficients, where for smooth \(f\) the functor \(f^!\) is a twist and a shift.

Let \(f_i:X_i\to Y_i\), \(i=1,2\), be morphisms of separated \(k\)-schemes of finite type. For \(M_i\in D^b_{\mathrm{ctf}}(X_i,A)\), using base change for compact supports and the projection formula, we obtain the Künneth isomorphism

\[
R(f_1\times f_2)_!(M_1\boxtimes M_2)
\simeq Rf_{1!}M_1\boxtimes Rf_{2!}M_2.
\]

For the structure maps it reads \(R\Gamma_c(X_1\times X_2,M_1\boxtimes M_2)\simeq R\Gamma_c(X_1,M_1)\otimes^{\mathbf L} R\Gamma_c(X_2,M_2)\). For \(L_i\in D^b_{\mathrm{ctf}}(Y_i,A)\) let

\[
\alpha:f_1^!L_1\boxtimes f_2^!L_2\longrightarrow
(f_1\times f_2)^!(L_1\boxtimes L_2)
\]

be the morphism adjoint, under \(R(f_1\times f_2)_!\dashv(f_1\times f_2)^!\), to the composite

\[
\begin{aligned}
R(f_1\times f_2)_!(f_1^!L_1\boxtimes f_2^!L_2)
&\simeq Rf_{1!}f_1^!L_1\boxtimes Rf_{2!}f_2^!L_2\\
&\xrightarrow{\varepsilon_{f_1}\boxtimes\varepsilon_{f_2}}
L_1\boxtimes L_2,
\end{aligned}
\]

where \(\varepsilon_{f_i}\) is the counit of \(Rf_{i!}\dashv f_i^!\).

% object: proposition-exceptional-external-kunneth (blueprint/chap-The-Lefschetz-Verdier-formula/sec-The-exceptional-Kunneth-formula/proposition-exceptional-external-kunneth/_section.md)
\phantomsection\label{bp19}%
\begin{bpstatementit}{Proposition~11.6}{ (Exceptional external Künneth)}
For all \(L_i\in D^b_{\mathrm{ctf}}(Y_i,A)\), the morphism \(\alpha\) is an isomorphism.
\end{bpstatementit}

% proof: proof/proposition-exceptional-external-kunneth
\begin{proof}[{Proof}]\phantomsection\label{bp20}%
Both sides of \(\alpha\) are exact in each variable, so by dévissage it is enough to check \(\alpha\) on generators. First, for constant coefficients we compute \(f_i^!\) by Theorem \bpout{9.3} for smooth maps and by localization, and check \(\alpha\) directly. Next, general coefficients are built from extensions by zero \(j_!F\) of lisse sheaves. Since \(f^!\) does not commute with \(j_!\), we replace \(j_!\) by \(Rj_*\), which does. The difference has smaller support, and we treat it by induction. For the classical statements we refer to [\hyperref[bp57]{SGA5}, III, (1.6.4), (1.7.3), Proposition 1.7.4 and (3.1.1)].

\emph{Constant and lisse coefficients.} First let the coefficients be constant. For smooth \(f_i\) of relative dimensions \(d_i\), Theorem \bpout{9.3} identifies \(f_i^!L_i\) with \(f_i^*L_i(d_i)[2d_i]\), and \((f_1\times f_2)^!\) with \((f_1\times f_2)^*(d_1+d_2)[2d_1+2d_2]\), so both sides of \(\alpha\) become the same twisted pullback. That \(\alpha\) is this identification follows from the compatibility of the trace map (Definition \bpout{9.24}) with products. Indeed, \(\alpha\) is defined through the counits, and for a smooth map the trace map is the counit composed with \(\gamma_f\).

For closed immersions we use the localization triangle and the case of ordinary external Künneth where one factor is the identity, applied to the complementary open immersions. In general, every map factors locally as a closed immersion followed by a smooth projection of affine spaces. Since \(\alpha\) is compatible with composition of morphisms and with restriction to open subschemes, Zariski localization gives the constant case for general maps. Next, lisse finite flat coefficients become constant with finite projective values on an étale cover, so descent gives the lisse case, and then the case of bounded perfect lisse models. In each of these steps the isomorphism we obtain is \(\alpha\) itself: for smooth maps because \(\gamma_f\) identifies the counit with the trace map, and for localization and descent because \(\alpha\) is natural and compatible with restriction.

\emph{General coefficients.} A constructible complex is lisse on a dense open, and the rest has smaller support. We therefore argue by induction on the sum of the dimensions of the supports of \(L_1,L_2\). Pushforward along closed immersions is compatible with \(\alpha\) by proper base change and the exceptional exchange of Corollary \bpout{10.9}. Hence we may first assume that the support is the entire base. As both sides of \(\alpha\) are exact in each variable, a finite constructible stratification, bounded flat models and two-out-of-three reduce us to generators \(j_{i!}F_i\) for dense open \(j_i:U_i\hookrightarrow Y_i\), with \(F_i\) a lisse sheaf with finite projective stalks (Lemma \bpout{3.13}).

Since \(f_i^!\) does not commute with \(j_{i!}\) in general, we pass through \(Rj_{i*}\). Let \(j'_i\) and \(f_{i,U}\) be the base changes of \(j_i\) and \(f_i\). The exceptional exchange \(f_i^!Rj_{i*}\simeq Rj'_{i*}f_{i,U}^!\) is the mate, under the right adjoints, of base change for \(Rf_{i!}\) along the open immersion. Hence it is compatible with the counits used to define \(\alpha\). By this exchange and the ordinary external Künneth formula, the comparison on \((Rj_{1*}F_1,Rj_{2*}F_2)\) is the pushforward of the lisse comparison. Let \(z_i:Z_i\hookrightarrow Y_i\) be the closed complement of \(U_i\). In the triangles

\[
j_{i!}F_i\longrightarrow Rj_{i*}F_i\longrightarrow
z_{i*}z_i^*Rj_{i*}F_i\xrightarrow{+1},
\]

the third term is still constructible of finite Tor amplitude by Proposition \bpout{7.1}, and its support has smaller dimension. Since \(\alpha\) is natural, it maps the corresponding triangles to each other in each variable. By induction and two-out-of-three, first in one variable and then in the other, we obtain \(\alpha\) for \((j_{1!}F_1,j_{2!}F_2)\).
\end{proof}

% text: prose-11-09
\phantomsection\label{bp21}%
Using the Künneth formulas for both kinds of direct image and for \(f^!\), we now identify cohomological correspondences with sections of \(c^!(K\boxtimes\mathbb D_XK)\).

% object: lemma-external-duality-trace-morphism (blueprint/chap-The-Lefschetz-Verdier-formula/sec-The-exceptional-Kunneth-formula/lemma-external-duality-trace-morphism/_section.md)
\phantomsection\label{bp22}%
\begin{bpstatementit}{Lemma~11.7}{ (External duality)}
Let \(K\in D^b_{\mathrm{ctf}}(X,A)\), let \(p_1,p_2:X\times X\to X\) be the projections, and put \(P_X=K\boxtimes \mathbb D_XK\in D^b_{\mathrm{ctf}}(X\times X,A)\). For a correspondence \(c:C\to X\times X\) there are canonical isomorphisms

\[
P_X\xrightarrow{\sim}\cHom_A(p_2^*K,p_1^!K),
\qquad
\mathcal H_c(K):=\cHom_A(c_2^*K,c_1^!K)
\xrightarrow{\sim}c^!P_X.
\]
\end{bpstatementit}

% proof: proof/lemma-external-duality-trace-morphism
\begin{proof}[{Proof}]\phantomsection\label{bp23}%
The exceptional Künneth formula (Proposition \hyperref[bp19]{11.6}) and the ordinary one give Künneth formulas for internal Hom and for Verdier duals, and these give the two displayed isomorphisms.

\emph{Internal Hom and Verdier duals.} On the point this is the identity \(\Hom(V_1,W_1)\otimes\Hom(V_2,W_2)=\Hom(V_1\otimes V_2,W_1\otimes W_2)\) for finite free modules. For \(E_i,L_i\in D^b_{\mathrm{ctf}}(X_i,A)\) the external comparison \(\cHom_A(E_1,L_1)\boxtimes\cHom_A(E_2,L_2)\to
\cHom_A(E_1\boxtimes E_2,L_1\boxtimes L_2)\) is an isomorphism. By dévissage it suffices to check this on generators \(E_i=Rh_{i!}F_i\), with \(h_i\) the inclusion of a stratum and \(F_i\) a lisse sheaf with finite projective stalks (Lemma \bpout{3.13}), for which

\[
\cHom_A(Rh_{i!}F_i,L_i)
\simeq Rh_{i*}(F_i^\vee\otimes^{\mathbf L} h_i^!L_i).
\]

By the ordinary Künneth formula and \(\alpha\), the exterior product of these expressions is identified with internal Hom on the product. Taking \(L_i=\omega_{X_i}\) and using \(\alpha\) for the structure maps, we get, for \(M_i\in D^b_{\mathrm{ctf}}(X_i,A)\),

\[
\mathbb D_{X_1\times X_2}(M_1\boxtimes M_2)
\simeq \mathbb D_{X_1}M_1\boxtimes \mathbb D_{X_2}M_2.
\]

\emph{The two isomorphisms.} The first isomorphism is the sheaf version of \(\Hom(V,W)\simeq W\otimes V^\vee\), and we obtain it by writing internal Hom as a dual of a tensor product. Indeed, for \(M\in D^b_{\mathrm{ctf}}(U,A)\), \(N\in D^b_{\mathrm{ctf}}(V,A)\) and the projections \(p_U,p_V\) of \(V\times U\), Verdier exchange and biduality let us rewrite \(\cHom_A(p_U^*M,p_V^!N)\) as \(\mathbb D_{V\times U}(\mathbb D_VN\boxtimes M)\). By the formula for duals of exterior products this is \(N\boxtimes \mathbb D_UM\). Taking \(M=N=K\), with the factors of the product ordered as above, we obtain the first isomorphism. The second isomorphism follows from the canonical exchange of exceptional inverse image with Hom, for \(P,Q\in D^b_{\mathrm{ctf}}(X\times X,A)\),

\[
c^!\cHom_A(P,Q)\simeq\cHom_A(c^*P,c^!Q),
\]

applied to \(P=p_2^*K\) and \(Q=p_1^!K\). To prove the exchange, we use Yoneda inside \(\Loc(D^b_{\mathrm{ctf}}(C,A))\), the category on which \(c_!\dashv c^!\) is defined (Theorem \bpout{5.2}). Test objects in all of \(D(C_{\mathrm{pro}\acute{e}t},A)\) would not do, since \(c^!\) is not defined there. Both sides of the exchange lie in \(D^b_{\mathrm{ctf}}(C,A)\) by Proposition \bpout{7.1}, so these test objects suffice. By Proposition \bpout{5.15} the localizing categories are \(D(C_{\acute{e}t},A)\) and \(D((X\times X)_{\acute{e}t},A)\), where the projection formula for \(c_!\) holds for all objects [Stacks, \href{https://stacks.math.columbia.edu/tag/0GL5}{0GL5}], and tensoring with \(c^*P\) preserves \(\Loc(D^b_{\mathrm{ctf}}(C,A))\). For \(M\) in \(\Loc(D^b_{\mathrm{ctf}}(C,A))\) we obtain

\[
\begin{aligned}
\Hom\bigl(M,c^!\cHom_A(P,Q)\bigr)
&\simeq\Hom\bigl(c_!M\otimes^{\mathbf L}P,Q\bigr)
\simeq\Hom\bigl(c_!(M\otimes^{\mathbf L}c^*P),Q\bigr)\\
&\simeq\Hom\bigl(M\otimes^{\mathbf L}c^*P,c^!Q\bigr)
\simeq\Hom\bigl(M,\cHom_A(c^*P,c^!Q)\bigr),
\end{aligned}
\]

by \(c_!\dashv c^!\) and the tensor–Hom adjunction, the projection formula, \(c_!\dashv c^!\) again, and the tensor–Hom adjunction.
\end{proof}

% text: prose-11-10
\phantomsection\label{bp24}%
In linear algebra the trace is evaluation \(W\otimes W^\vee\to A\) on the diagonal. Here evaluation takes values in \(\omega_X\) on the diagonal, and by applying \(c^!\) we move it to the fixed locus.

% object: definition-trace-morphism (blueprint/chap-The-Lefschetz-Verdier-formula/sec-The-exceptional-Kunneth-formula/definition-trace-morphism/_section.md)
\phantomsection\label{bp25}%
\begin{bpstatement}{Definition~11.8}{ (Trace morphism)}
In the setting of Lemma \hyperref[bp22]{11.7}, let \(e_K:P_X\to R\Delta_{X*}\omega_X\) be adjoint to the evaluation \(\Delta_X^*P_X=K\otimes^{\mathbf L} \mathbb D_XK\to\omega_X\). The \emph{trace morphism} is the composite

\[
\mathcal H_c(K)\simeq c^!P_X\xrightarrow{c^!e_K}c^!R\Delta_{X*}\omega_X
\xrightarrow{\sim}Ri_{c*}d_c^!\omega_X
\simeq Ri_{c*}\omega_{\Fix(c)}
\]

of the second isomorphism of the lemma, \(c^!e_K\), exceptional base change for the proper map \(\Delta_X\), and the composition of exceptional inverse images \(d_c^!a_X^!=a_{\Fix(c)}^!\).
\end{bpstatement}


% object: sec-Trace-classes-and-local-terms (blueprint/chap-The-Lefschetz-Verdier-formula/sec-Trace-classes-and-local-terms/_section.md), content hidden
\subsection*{\texorpdfstring{11.4\quad Trace classes and local terms}{11.4 Trace classes and local terms}}
\phantomsection\label{bp26}%
\addcontentsline{toc}{subsection}{\texorpdfstring{11.4\quad Trace classes and local terms}{11.4 Trace classes and local terms}}
% text: prose-11-11
\phantomsection\label{bp27}%
We want to write \(\operatorname{tr}(T_c(u))\) as a sum of contributions from the fixed locus. The first guess, a sum over fixed points of the traces of \(u\) on stalks, does not make sense in general, because the fixed locus need not be finite. For example, for the identity correspondence of \(\mathbf P^1\) and \(K=A\), every point is fixed. The trace on \(R\Gamma(\mathbf P^1,A)=A\oplus A(-1)[-2]\) is \(1+1=2\), the self-intersection of the diagonal.

What does make sense is a class on \(\Fix(c)\). By Lemma \hyperref[bp22]{11.7}, \(u\) is a section of \(c^!(K\boxtimes\mathbb D_XK)\), and composing with the trace morphism we obtain a section of \(\omega_{\Fix(c)}\). Such a section can be integrated over a proper piece \(B\) of the fixed locus. Indeed, \(H^0(B,\omega_B)=\operatorname{Hom}(Ra_{B!}A_B,A)\), and \(Ra_{B*}A_B=Ra_{B!}A_B\) for proper \(B\). Composing with the unit \(A\to Ra_{B*}A_B\) therefore gives an element of \(A\). In the example above, \(H^0(\mathbf P^1,\omega_{\mathbf P^1})=H^2(\mathbf P^1,A(1))\) and the degree map \(H^2(\mathbf P^1,A(1))\to A\) is an isomorphism, so Theorem \hyperref[bp37]{11.13} shows that the trace class is \(2\) times the class of a point.

We will see that the trace class is additive in \(K\) (Proposition \hyperref[bp54]{11.19}) and compatible with proper pushforward (Theorem \hyperref[bp37]{11.13}). However, at an isolated fixed point the local term can still differ from the trace on the stalk (Remark \bpout{12.4}). For Frobenius the two are shown to agree in Proposition \bpout{12.16}.

% object: definition-trace-class-local-terms (blueprint/chap-The-Lefschetz-Verdier-formula/sec-Trace-classes-and-local-terms/definition-trace-class-local-terms/_section.md)
\phantomsection\label{bp28}%
\begin{bpstatement}{Definition~11.9}{ (Trace class and local terms)}
Let \(u:c_2^*K\to c_1^!K\) be a cohomological correspondence on \((X,K)\). By adjunction, \(u\) is a section \(A_C\to\mathcal H_c(K)\). Its composite with the trace morphism of Definition \hyperref[bp25]{11.8}, transposed by \(i_c^*\dashv Ri_{c*}\), is the \emph{trace class}

\[
\operatorname{Tr}_c(u):A_{\Fix(c)}\longrightarrow\omega_{\Fix(c)}.
\]

For a proper \(k\)-scheme \(B\), integration is the counit map

\[
\int_B:H^0(B,\omega_B)
=H^0(Ra_{B!}a_B^!A_s)\longrightarrow A.
\]

If \(B\subset \Fix(c)\) is proper and open and closed, the \emph{local term} of \(u\) on \(B\) is

\[
\operatorname{LT}_B(u)=\int_B\operatorname{Tr}_c(u)|_B.
\]
\end{bpstatement}

% text: prose-11-12
\phantomsection\label{bp29}%
On the point, the local term is the categorical trace of Definition \hyperref[bp5]{11.1}, and we use this later to identify the local terms of Frobenius with traces on stalks.

% object: lemma-local-terms-point (blueprint/chap-The-Lefschetz-Verdier-formula/sec-Trace-classes-and-local-terms/lemma-local-terms-point/_section.md)
\phantomsection\label{bp30}%
\begin{bpstatementit}{Lemma~11.10}{ (Local terms on the point)}
Let \(P\in\operatorname{Perf}(A)\), viewed as an object of \(D^b_{\mathrm{ctf}}(s,A)\), and let \(u:P\to P\) be a cohomological correspondence on \((s,P)\) for the identity correspondence of \(s\), whose fixed locus is \(s\). Then

\[
\operatorname{LT}_s(u)=\operatorname{tr}_A(u\mid P).
\]
\end{bpstatementit}

% proof: proof/lemma-local-terms-point
\begin{proof}[{Proof}]\phantomsection\label{bp31}%
On the point each step of the construction of the trace class is a step of the categorical trace. Here \(\omega_s=A_s\), so \(\mathbb D_sP=P^\vee\), and \(P_s=P\boxtimes
P^\vee=P\otimes^{\mathbf L}P^\vee\). First, under the first isomorphism of Lemma \hyperref[bp22]{11.7}, the map \(u\) becomes the element of \(P\otimes^{\mathbf L}P^\vee\) obtained by applying \(u\otimes1\) to the coevaluation. Next, the diagonal of \(s\) is the identity, so the trace morphism is the evaluation \(P\otimes^{\mathbf L}P^\vee\to A\), taken after the symmetry that puts the factors in the order \(P^\vee\otimes P\). That symmetry is precisely the symmetry in Definition \hyperref[bp5]{11.1}, and integration over \(s\) is the identity of \(A\). The composite is therefore the categorical trace.
\end{proof}

% text: prose-11-13
\phantomsection\label{bp32}%
The global side of the trace formula is the trace of the following endomorphism.

% object: definition-action-correspondence (blueprint/chap-The-Lefschetz-Verdier-formula/sec-Trace-classes-and-local-terms/definition-action-correspondence/_section.md)
\phantomsection\label{bp33}%
\begin{bpstatement}{Definition~11.11}{ (Action of a correspondence on cohomology)}
Let \(u\) be a cohomological correspondence on \((X,K)\) for a correspondence \(c\) with \(c_2\) proper. The endomorphism \(T_c(u)\) of \(R\Gamma_c(X,K)\) is the composite

\[
\begin{aligned}
T_c(u):R\Gamma_c(X,K)&\xrightarrow{\eta_{c_2}}
R\Gamma_c(C,c_2^*K)\xrightarrow{u}R\Gamma_c(C,c_1^!K)\\
&\simeq R\Gamma_c(X,Rc_{1!}c_1^!K)
\xrightarrow{\varepsilon_{c_1}}R\Gamma_c(X,K).
\end{aligned}
\]
\end{bpstatement}

% text: prose-11-14
\phantomsection\label{bp34}%
In the first arrow we use the ordinary unit together with \(Rc_{2*}=Rc_{2!}\), which holds since \(c_2\) is proper, and the displayed isomorphism comes from the composition of direct images with compact support.

The Lefschetz–Verdier formula computes the trace of \(T_c(u)\) in terms of the trace class. We obtain it as the special case, for the map to a point, of the compatibility of trace classes with proper pushforward. This works because \(T_c(u)\) is itself a cohomological correspondence: for proper \(X\) and \(C\) it is a correspondence on the point with coefficients \(R\Gamma(X,K)\), and on the point the trace class is the trace by Lemma \hyperref[bp30]{11.10}. The global trace is therefore the trace class of a pushed-forward correspondence, and it remains to show that trace classes commute with pushforward along proper maps.

To push a cohomological correspondence forward along proper morphisms, we must produce from \(u\) a map \(c_2'^*Rf_*K\to c_1'^!Rf_*K\). The only maps available for this are the units and counits of the adjunctions on the two legs, used as in the construction of \(T_c(u)\). Here properness of \(f\) and \(g\) gives \(Rf_*=Rf_!\) and \(Rg_*=Rg_!\), and this allows us to push forward the counit for \(Rc_{1!}\dashv c_1^!\). We also push forward along closed immersions (Section \hyperref[bp40]{11.5}).

% object: definition-pushforward-correspondence (blueprint/chap-The-Lefschetz-Verdier-formula/sec-Trace-classes-and-local-terms/definition-pushforward-correspondence/_section.md)
\phantomsection\label{bp35}%
\begin{bpstatement}{Definition~11.12}{ (Pushforward of a correspondence)}
Let \(f:X\to Y\) and \(g:C\to C'\) be proper morphisms, and let \(c'=(c'_1,c'_2):C'\to Y\times Y\) be a correspondence with \(c'_ig=fc_i\) for \(i=1,2\). Put \(K'=Rf_*K\in D^b_{\mathrm{ctf}}(Y,A)\). For a cohomological correspondence \(u\) on \((X,K)\), the \emph{pushforward} \(u':c_2'^*K'\to c_1'^!K'\) of \(u\) is the composite \(u'=\beta\,Rg_*u\,\alpha\), where \(\alpha\) is built from the unit and counit of \((\,)^*\dashv R(\,)_*\),

\[
\alpha:c_2'^*Rf_*K\xrightarrow{\eta_g}
Rg_*c_2^*f^*Rf_*K\xrightarrow{\varepsilon_f}Rg_*c_2^*K,
\]

and \(\beta:Rg_*c_1^!K\to c_1'^!Rf_*K\) is built from the counit of \(R(\,)_!\dashv(\,)^!\): it is adjoint to

\[
Rc'_{1!}Rg_*c_1^!K\simeq Rf_*Rc_{1!}c_1^!K
\xrightarrow{Rf_*\varepsilon_{c_1}}Rf_*K.
\]

The map \(g\) induces a proper morphism \(h:\Fix(c)\to \Fix(c')\) of fixed loci, and \(h_*:H^0(\Fix(c),\omega_{\Fix(c)})\to H^0(\Fix(c'),\omega_{\Fix(c')})\) denotes the map given by the ordinary unit, proper pushforward and the exceptional counit.
\end{bpstatement}

% text: prose-11-15
\phantomsection\label{bp36}%
For \(Y=s\), \(C'=s\) and \(f=a_X\), \(g=a_C\) with \(X\) and \(C\) proper, the pushforward \(u'\) is an endomorphism of \(R\Gamma(X,K)\). Unwinding \(\alpha\) and \(\beta\) shows that it is \(T_c(u)\), and the following theorem compares trace classes before and after pushforward. In this special case it is the Lefschetz–Verdier formula.

% object: theorem-proper-pushforward-lefschetz-verdier (blueprint/chap-The-Lefschetz-Verdier-formula/sec-Trace-classes-and-local-terms/theorem-proper-pushforward-lefschetz-verdier/_section.md)
\phantomsection\label{bp37}%
\begin{bpstatementit}{Theorem~11.13}{ (Proper pushforward and Lefschetz–Verdier)}
Let \(f:X\to Y\) and \(g:C\to C'\) be proper, and let \(c'=(c'_1,c'_2):C'\to Y\times Y\) satisfy \(c'_ig=fc_i\). Set \(K'=Rf_*K\), and let \(u'\) be the pushforward of \(u\) to \((Y,K')\) of Definition \hyperref[bp35]{11.12}. With the induced proper map \(h:\Fix(c)\to \Fix(c')\) one has

\[
\operatorname{Tr}_{c'}(u')=h_*\operatorname{Tr}_c(u).
\]

If \(X\) and \(C\) are proper over \(k\), this gives

\[
\operatorname{tr}_A(T_c(u)\mid R\Gamma(X,K))
=\int_{\Fix(c)}\operatorname{Tr}_c(u)
=\sum_B\operatorname{LT}_B(u),
\]

for any finite open-and-closed decomposition \(\Fix(c)=\coprod B\).
\end{bpstatementit}

% proof: proof/theorem-proper-pushforward-lefschetz-verdier
\begin{proof}[{Proof}]\phantomsection\label{bp38}%
The trace class of \(u\) is built in three steps: \(u\) becomes a map into \(P_X=K\boxtimes\mathbb D_XK\), evaluation maps \(P_X\) to the dualizing complex on the diagonal, and \(c^!\) moves the result to the fixed locus. We check that each step commutes with pushforward. For the first two steps this means identifying the map adjoint to \(u'\) with the pushforward of the map adjoint to \(u\), and evaluation for \(K'\) with the pushforward of evaluation for \(K\). The third step needs more care. The squares \(c'_ig=fc_i\) are only required to commute, so in general \(\Fix(c)\ne C\times_{C'}\Fix(c')\), and we compare the two fixed loci through an intermediate scheme.

\emph{Maps into \(P\).} Put \(P_Y=K'\boxtimes \mathbb D_YK'\), and let \(\tau_f:Rf_*\omega_X\to\omega_Y\) be the counit. By the Künneth formula for the proper map \(f\times f\) and the exchange \(\mathbb D_YRf_*\simeq Rf_*\mathbb D_X\) of Verdier duality, we obtain an isomorphism

\[
\rho_f:R(f\times f)_*P_X\xrightarrow{\sim}P_Y.
\]

A cohomological correspondence \(u\) defines a map \(\theta_u:Rc_!A_C\to P_X\), obtained by exceptional adjunction from the section \(A_C\to c^!P_X\). Under \(Rc'_!Rg_*\simeq R(f\times f)_*Rc_!\), the map defined by the pushforward \(u'\) is

\[
\theta_{u'}=
\rho_f\,R(f\times f)_*\theta_u\,
(Rc'_!Rg_*\simeq R(f\times f)_*Rc_!)\,Rc'_!\eta_g.
\]

To verify this identity, we apply tensor–Hom adjunction and then the exceptional adjunction for \(c'_1\). After this both sides become \(Rc'_{1!}\alpha\), followed by \(Rf_*Rc_{1!}u\) and \(Rf_*\varepsilon_{c_1}\). The two sides use the projection formulas for \(c_1\) and for \(c'_1\) in different orders, but they agree because \(c'_1g=fc_1\) and projection formulas are compatible with composition.

\emph{Evaluation.} Evaluation is compatible with \(\rho_f\), in the sense that

\[
e_{K'}\rho_f=
R\Delta_{Y*}\tau_f\,R(f\times f)_*e_K,
\qquad (f\times f)\Delta_X=\Delta_Yf.
\]

After restricting to the diagonal and taking tensor–Hom adjoints, this says that Verdier duality for the proper map \(f\) identifies evaluation against \(Rf_*K\) with the pushforward of evaluation against \(K\), followed by \(\tau_f\). This holds because that duality is defined by means of the counit \(\tau_f\).

\emph{Fixed loci.} The two maps into \(P_Y\) and the two evaluations therefore match. The trace classes, however, live on two different fixed loci, and we must see that passing from \(\Fix(c)\) to \(\Fix(c')\) is the map \(h_*\) of Definition \hyperref[bp35]{11.12}. Recall that the fixed loci are the pullbacks of the diagonals, \(\Fix(c)=C\times_{X\times X}X\) and \(\Fix(c')=C'\times_{Y\times Y}Y\). To compare them, we factor \(h\) through the fibre product \(C\times_{C'}\Fix(c')\):

\[
\Fix(c)\xrightarrow{i'}C\times_{C'}\Fix(c')\xrightarrow{\pi}\Fix(c'),
\]

where the closed immersion \(i'\) is the pullback of \(X\to X\times_YX\) and \(\pi\) is the base change of \(g\). We use proper base change for the square defining \(C\times_{C'}\Fix(c')\), and the exceptional counit for \(i'\). Since counits are compatible with composition, the integration map for \(h\) is \(\tau_h=\tau_\pi R\pi_*\tau_{i'}\). Applying the identity for \(\theta_{u'}\) and the compatibility of evaluation, together with these base change maps, we identify the map \(A_{C'}\to Ri_{c'*}\omega_{\Fix(c')}\) defining the left-hand trace with

\[
A_{C'}\xrightarrow{\eta_g}Rg_*A_C
\longrightarrow Rg_*Ri_{c*}\omega_{\Fix(c)}
=Ri_{c'*}Rh_*\omega_{\Fix(c)}
\xrightarrow{Ri_{c'*}\tau_h}Ri_{c'*}\omega_{\Fix(c')}.
\]

Taking the \(i_{c'}^*\)-adjoint by means of the identity of composite units \(Rg_*\eta_{i_c}\,\eta_g=Ri_{c'*}\eta_h\,\eta_{i_{c'}}\), we obtain the first formula (compare [\hyperref[bp57]{SGA5}, III, Theorem 4.4 and Corollary 4.7]).

\emph{The trace formula.} For proper \(X\) and \(C\), we push forward to the identity correspondence on \(s\), using \(f=a_X\) and \(g=a_C\). Then the pushed cohomological correspondence is \(T_c(u)\). Indeed, for \(g=a_C\) the map \(\alpha\) is the pullback \(R\Gamma(X,K)\to R\Gamma(C,c_2^*K)\) and \(\beta\) is the counit of \(c_1\) after \(R\Gamma\), and these are the outer arrows in the definition of \(T_c(u)\). Moreover, the fixed locus of the identity correspondence of \(s\) is \(s\) itself. Hence the left side of the first formula is the local term of \(T_c(u)\) on \(s\), which by Lemma \hyperref[bp30]{11.10} is the trace of \(T_c(u)\). On the right, \(h_*\) is integration over \(\Fix(c)\), and since integration over a disjoint union is the sum of the integrations, this gives the second assertion.
\end{proof}

% object: example-correspondence-on-point (blueprint/chap-The-Lefschetz-Verdier-formula/sec-Trace-classes-and-local-terms/example-correspondence-on-point/_section.md)
\phantomsection\label{bp39}%
\begin{bpstatement}{Example~11.14}{ (A correspondence on the point)}
Let \(X=s\), let \(C=\coprod_{i=1}^ns\) be \(n\) points, with \(c_1=c_2\) the structure map, and let \(u:c_2^*A=A_C\to c_1^!A=A_C\) be the identity, using \(c_1^!=c_1^*\) for the finite étale map \(c_1\). Then pullback along \(c_2\) is the diagonal \(A\to A^n\), and the counit \(Rc_{1!}c_1^!A=A^n\to A\) is the sum. Hence \(T_c(u)\) is multiplication by \(n\) on \(R\Gamma(s,A)=A\), with trace \(n\). Since \(X\times X=s\), the fixed locus is all of \(C\). On each of these \(n\) points \(u\) restricts to the identity correspondence of \(s\) on \(A\), whose local term is \(\operatorname{tr}_A(1\mid A)=1\) by Lemma \hyperref[bp30]{11.10}. The right side of the theorem is therefore \(n\) as well. Note that the local terms are indexed by the points of \(C\), not of \(X\): here \(X\) has one point and the fixed locus has \(n\).
\end{bpstatement}


% object: sec-Additivity-of-trace-classes (blueprint/chap-The-Lefschetz-Verdier-formula/sec-Additivity-of-trace-classes/_section.md), content hidden
\subsection*{\texorpdfstring{11.5\quad Additivity of trace classes}{11.5 Additivity of trace classes}}
\phantomsection\label{bp40}%
\addcontentsline{toc}{subsection}{\texorpdfstring{11.5\quad Additivity of trace classes}{11.5 Additivity of trace classes}}
% text: prose-11-16
\phantomsection\label{bp41}%
To compute local terms, we cut \(K\) into pieces along a decomposition of \(X\) into an open and a closed part, both invariant under the correspondence. For this the trace class must be additive. For the trace of an endomorphism of a perfect complex, additivity says that a block triangular matrix has trace equal to the sum of the traces of its diagonal blocks. By Remark \hyperref[bp10]{11.3} it can only be expected for filtrations lifted to the filtered derived category. The trace class needs the same care, and we lift each step of the construction of the trace morphism to filtered complexes.

In a two-step filtration, for example, a filtered correspondence is a block triangular matrix \(\left(\begin{smallmatrix}u_0&*\\0&u_1\end{smallmatrix}\right)\) whose off-diagonal entry has weight \(1\). Since \(\omega_X\) sits in filtration degree \(0\), the trace class does not see this entry (Lemma \hyperref[bp45]{11.16} below).

The operations that enter the trace morphism have filtered versions, with explicit formulas for their graded pieces. We use finite decreasing filtrations, the convolution filtration on tensor products, and the filtration on internal Hom given by

\[
F^r\mathcal Hom(\mathcal M,\mathcal N)
=\{v:v(F^a\mathcal M)\subset F^{a+r}\mathcal N
\text{ for every }a\}.
\]

Write \(\cHom_F\), \(\otimes_F\) and \(\mathbb D_F=\cHom_F(-,\omega_X)\) for the derived functors on the filtered derived category, where \(\omega_X\) carries the filtration concentrated in weight zero, and \(M_a\), \(N_a\), \(K_a\) for the graded pieces of \(\mathcal M\), \(\mathcal N\), \(\mathcal K\).

% object: lemma-graded-pieces-filtered-operations (blueprint/chap-The-Lefschetz-Verdier-formula/sec-Additivity-of-trace-classes/lemma-graded-pieces-filtered-operations/_section.md)
\phantomsection\label{bp42}%
\begin{bpstatementit}{Lemma~11.15}{ (Graded pieces of filtered operations)}
For finitely filtered complexes \(\mathcal M,\mathcal N,\mathcal K\) on \(X\) with graded pieces in \(D^b_{\mathrm{ctf}}(X,A)\) there are isomorphisms

\[
\begin{aligned}
\operatorname{gr}^r\cHom_F(\mathcal M,\mathcal N)
&\simeq\bigoplus_a\cHom(M_a,N_{a+r}),\\
\operatorname{gr}^r\mathbb D_F\mathcal K&\simeq \mathbb D_XK_{-r},\\
\operatorname{gr}^r(\mathcal M\otimes_F\mathcal N)
&\simeq\bigoplus_{a+b=r}M_a\otimes^{\mathbf L} N_b.
\end{aligned}
\]

Under these isomorphisms, the restriction of the filtered evaluation \(\mathcal K\otimes_F\mathbb D_F\mathcal K\to\omega_X\) to the summand \(K_a\otimes^{\mathbf L}\mathbb D_XK_a\) of \(\operatorname{gr}^0\) is the evaluation of \(K_a\).
\end{bpstatementit}

% proof: proof/lemma-graded-pieces-filtered-operations
\begin{proof}[{Proof}]\phantomsection\label{bp43}%
The operations are derived using filtered flat and filtered injective resolutions. For finite filtrations these resolutions can be built successively on the graded pieces by mapping cylinders, and on them the three formulas hold termwise. For the evaluation, recall that the filtered evaluation is the tensor–Hom adjoint of \(\operatorname{id}_{\mathcal K}\), whose graded map in weight zero is \(\operatorname{id}_{K_a}\) on each summand. Hence on \(K_a\otimes^{\mathbf L}\mathbb D_XK_a\) it is the ordinary evaluation. Moreover, the Koszul signs are those of the underlying complexes, and the filtration weights contribute no signs.
\end{proof}

% text: prose-11-17
\phantomsection\label{bp44}%
With these formulas we lift the trace morphism of Definition \hyperref[bp25]{11.8} to filtered complexes.

% object: lemma-filtered-trace-morphism (blueprint/chap-The-Lefschetz-Verdier-formula/sec-Additivity-of-trace-classes/lemma-filtered-trace-morphism/_section.md)
\phantomsection\label{bp45}%
\begin{bpstatementit}{Lemma~11.16}{ (Filtered trace morphism)}
Let \(c:C\to X\times X\) be a correspondence and let \(\mathcal K\) be a finitely filtered complex on \(X\) with underlying complex \(K\) and graded pieces \(K_a\) in \(D^b_{\mathrm{ctf}}(X,A)\). There is a morphism of filtered complexes

\[
\operatorname{tr}_F:
\cHom_F(c_2^*\mathcal K,c_1^!\mathcal K)\longrightarrow
Ri_{c*}\omega_{\Fix(c)},
\]

with the target in weight zero, whose underlying morphism is the trace morphism for \(K\), and whose restriction to \(F^0\) factors through

\[
\operatorname{gr}^0\cHom_F(c_2^*\mathcal K,c_1^!\mathcal K)
\simeq\bigoplus_a\cHom(c_2^*K_a,c_1^!K_a)
\]

by a map whose restriction to the summand \(a\) is the trace morphism for \(K_a\).
\end{bpstatementit}

% proof: proof/lemma-filtered-trace-morphism
\begin{proof}[{Proof}]\phantomsection\label{bp46}%
The trace morphism is the composite of four steps: the exceptional tensor–Hom adjunction, external duality, the adjoint of evaluation along the diagonal, and exceptional proper base change for the closed diagonal. Since each is defined by units, counits and tensor evaluation, the composite has a filtered version

\[
\cHom_F(c_2^*\mathcal K,c_1^!\mathcal K)
\simeq c^!(\mathcal K\boxtimes_F\mathbb D_F\mathcal K)
\longrightarrow Ri_{c*}\omega_{\Fix(c)}
\]

whose graded pieces, computed by Lemma \hyperref[bp42]{11.15}, are the unfiltered versions. In particular the filtered comparison for external duality is an isomorphism, because each of its graded maps is a finite sum of the isomorphisms in Lemma \hyperref[bp22]{11.7}, and graded pieces detect isomorphisms of finitely filtered complexes. Forgetting the filtration we recover the trace morphism for \(K\).

The target is pure of weight zero, so \(\operatorname{tr}_F\) vanishes on \(F^1\), and its restriction to \(F^0\) factors through \(\operatorname{gr}^0\). This factorization holds on filtered dg models, hence also in the filtered derived category. By Lemma \hyperref[bp42]{11.15}, \(\operatorname{gr}^0\) of the source is \(\bigoplus_a\cHom(c_2^*K_a,c_1^!K_a)\). Now the steps for internal Hom and external duality preserve the two indices, and on a pair of equal indices evaluation along the diagonal is the evaluation of \(K_a\). Hence the map on the summand \(a\) is the trace morphism for \(K_a\).
\end{proof}

% text: prose-11-18
\phantomsection\label{bp47}%
To apply this to an open subscheme \(U\) and its closed complement \(Z\), we must lift \(u\) to the localization triangle. For this the correspondence must respect the decomposition: we need \(c_1^{-1}(Z)\subset c_2^{-1}(Z)\). For the graph of an endomorphism \(\varphi\), where \(c_1=\operatorname{id}_X\) and \(c_2=\varphi\), this says that \(\varphi(Z)\subset Z\).

It is convenient to work with exceptional adjoints. Put \(\Theta=Rc_{1!}c_2^*:D^b_{\mathrm{ctf}}(X,A)\to D^b_{\mathrm{ctf}}(X,A)\). In these terms a cohomological correspondence \(u\) on \((X,K)\) is the same as a map \(\widehat u:\Theta K\to K\), its adjoint under \(Rc_{1!}\dashv c_1^!\). Under the invariance condition there are no nonzero maps from \(\Theta K_U\) to \(K_Z\), so \(u\) restricts to \(K_U\) in an essentially unique way. Let

\[
C_U=c_2^{-1}(U),\qquad C_Z=(c_1^{-1}(Z))_{\mathrm{red}},\qquad
i':C_Z\longrightarrow C,
\]

with \(i'\) the closed immersion. By the invariance condition both legs of \(C_U\) land in \(U\) and both legs of \(C_Z\) land in \(Z\); we write \(c_U=(c_{1U},c_{2U}):C_U\to U\times U\) and \(c_Z=(c_{1Z},c_{2Z}):C_Z\to Z\times Z\) for the restricted correspondences.

% object: lemma-lift-correspondence-invariant-decomposition (blueprint/chap-The-Lefschetz-Verdier-formula/sec-Additivity-of-trace-classes/lemma-lift-correspondence-invariant-decomposition/_section.md)
\phantomsection\label{bp48}%
\begin{bpstatementit}{Lemma~11.17}{ (Lifting a correspondence to an invariant decomposition)}
Let \(c:C\to X\times X\) be a correspondence, \(i:Z\to X\) a closed immersion with open complement \(j:U\to X\), and suppose

\[
c_1^{-1}(Z)\subset c_2^{-1}(Z),
\qquad\text{equivalently }c_2^{-1}(U)\subset c_1^{-1}(U).
\]

Let \(u\) be a cohomological correspondence on \((X,K)\), and write \(K_U=j_!j^*K\) and \(K_Z=Ri_*i^*K\) for the two pieces of the localization triangle \(K_U\to K\to K_Z\). Then \(u\) has a lift, unique up to contractible choice, to a cohomological correspondence \(\tilde u\) in the filtered derived category on the two-step filtered complex given by this triangle. Its graded pieces are cohomological correspondences \(u_U\) on \((X,K_U)\) and \(u_Z\) on \((X,K_Z)\), and \(\widehat u_Z\) restricted to \(Z\) is

\[
Rc_{1Z!}c_{2Z}^*i^*K\simeq i^*\Theta K
\xrightarrow{i^*\widehat u}i^*K.
\]
\end{bpstatementit}

% proof: proof/lemma-lift-correspondence-invariant-decomposition
\begin{proof}[{Proof}]\phantomsection\label{bp49}%
By the hypothesis \(c_1^{-1}(Z)\subset c_2^{-1}(Z)\), base change of \(Rc_{1!}\) along \(i\), and the adjunction \(i^*\dashv Ri_*\), we have

\[
i^*\Theta K_U=0,\qquad
R\operatorname{Hom}_X(\Theta K_U,K_Z)
\simeq R\operatorname{Hom}_Z(i^*\Theta K_U,i^*K)=0.
\]

The second vanishing is that of the whole mapping complex, including all shifts. Hence the composite \(\Theta K_U\to\Theta K\xrightarrow{\widehat u}K\) factors through \(K_U\to K\), with a contractible space of choices. The coherent cofibre of this factorization is a morphism of two-step filtered objects

\[
\begin{array}{ccccc}
\Theta K_U&\longrightarrow&\Theta K&\longrightarrow&\Theta K_Z\\
\downarrow&&\downarrow \widehat u&&\downarrow\\
K_U&\longrightarrow&K&\longrightarrow&K_Z.
\end{array}
\]

Since this cofibre includes the homotopy, it is a morphism in the filtered derived category, and its adjoint under the filtered adjunction \(Rc_{1!}\dashv c_1^!\) is the lift \(\tilde u\). By base change, composition and nilpotent invariance (which lets us use the reduced scheme \(C_Z\)),

\[
\Theta K_U\simeq j_!Rc_{1U!}c_{2U}^*j^*K,
\qquad i^*\Theta K\simeq Rc_{1Z!}c_{2Z}^*i^*K.
\]

Applying \(i^*\) to the right-hand square and using \(i^*\Theta K_U=0\), so that \(i^*\Theta K\simeq i^*\Theta K_Z\), we obtain the formula for \(\widehat u_Z\) restricted to \(Z\).
\end{proof}

% text: prose-11-19
\phantomsection\label{bp50}%
The graded piece on the closed part is a correspondence on \(Z\) pushed forward to \(X\). This is the form in which Proposition \bpout{12.16} uses it, combined with Theorem \hyperref[bp37]{11.13}.

% object: lemma-closed-piece-pushforward (blueprint/chap-The-Lefschetz-Verdier-formula/sec-Additivity-of-trace-classes/lemma-closed-piece-pushforward/_section.md)
\phantomsection\label{bp51}%
\begin{bpstatementit}{Lemma~11.18}{ (The closed piece as a pushforward)}
In the situation of Lemma \hyperref[bp48]{11.17}, put \(M=i^*K\in D^b_{\mathrm{ctf}}(Z,A)\), and let \(u_{|Z}\) be the cohomological correspondence on \((Z,M)\) for \(c_Z\) whose exceptional adjoint is \(Rc_{1Z!}c_{2Z}^*M\simeq i^*\Theta K\xrightarrow{i^*\widehat u}M\). Then \(u_Z\) is the pushforward of \(u_{|Z}\) along the proper maps \(i\) and \(i'\) (Definition \hyperref[bp35]{11.12}).
\end{bpstatementit}

% proof: proof/lemma-closed-piece-pushforward
\begin{proof}[{Proof}]\phantomsection\label{bp52}%
The maps \(\alpha\) and \(\beta\) of Definition \hyperref[bp35]{11.12} are here

\[
c_2^*Ri_*M\xrightarrow{\eta_{i'}}
Ri'_*i'^*c_2^*Ri_*M
=Ri'_*c_{2Z}^*i^*Ri_*M
\xrightarrow{Ri'_*c_{2Z}^*\varepsilon_i}Ri'_*c_{2Z}^*M
\]

and the map \(Ri'_*c_{1Z}^!M\to c_1^!Ri_*M\) adjoint to

\[
Rc_{1!}Ri'_*c_{1Z}^!M
=Ri_*Rc_{1Z!}c_{1Z}^!M\xrightarrow{Ri_*\varepsilon_{c_{1Z}}}Ri_*M.
\]

The exceptional adjoint of the pushforward is thus a map \(\Theta Ri_*M\to Ri_*M\). After applying \(i^*\) and the base change identity \(i^*\Theta Ri_*M\simeq Rc_{1Z!}c_{2Z}^*M\), it becomes the exceptional adjoint of \(u_{|Z}\). By Lemma \hyperref[bp48]{11.17} the same holds for \(\widehat u_Z\). Since \(K_Z=Ri_*M\) and a map into \(Ri_*M\) is determined by its adjoint under \(i^*\dashv Ri_*\), the two correspondences agree.
\end{proof}

% text: prose-11-20
\phantomsection\label{bp53}%
Additivity takes the following form; in Proposition \bpout{12.16} we apply it with \(Z\) a single fixed point of Frobenius.

% object: proposition-additivity-trace-classes (blueprint/chap-The-Lefschetz-Verdier-formula/sec-Additivity-of-trace-classes/proposition-additivity-trace-classes/_section.md)
\phantomsection\label{bp54}%
\begin{bpstatementit}{Proposition~11.19}{ (Additivity of trace classes)}
Let \(c:C\to X\times X\) be a correspondence.

\begin{enumerate}
\item Let \(\mathcal K\) be a finitely filtered complex on \(X\) with underlying complex \(K\) and graded pieces \(K_a\) in \(D^b_{\mathrm{ctf}}(X,A)\), and let \(\tilde u:c_2^*\mathcal K\to c_1^!\mathcal K\) be a cohomological correspondence in the filtered derived category, with underlying correspondence \(u\) on \(K\) and graded pieces \(u_a\) on \(K_a\). Then \[
   \operatorname{Tr}_c(u)=\sum_a\operatorname{Tr}_c(u_a)
   \quad\text{in }H^0(\Fix(c),\omega_{\Fix(c)}).
   \]
\item Let \(i:Z\to X\) be a closed immersion with open complement \(U\) such that \(c_1^{-1}(Z)\subset c_2^{-1}(Z)\), let \(K\in D^b_{\mathrm{ctf}}(X,A)\), and let \(u\) be a cohomological correspondence on \((X,K)\). Then \[
   \operatorname{Tr}_c(u)=\operatorname{Tr}_c(u_U)+\operatorname{Tr}_c(u_Z),
   \] where \(u_U\) and \(u_Z\) are the graded pieces of the lift of Lemma \hyperref[bp48]{11.17}, and \(u_Z\) is the pushforward along \(i\) and \(C_Z\to C\) of the correspondence \(u_{|Z}\) on \((Z,i^*K)\) for \(c_Z:C_Z=(c_1^{-1}(Z))_{\mathrm{red}}\to
Z\times Z\) (Lemma \hyperref[bp51]{11.18}).
\end{enumerate}
\end{bpstatementit}

% proof: proof/proposition-additivity-trace-classes
\begin{proof}[{Proof}]\phantomsection\label{bp55}%
For (1), the section \(A_C\to\cHom_F(c_2^*\mathcal K,c_1^!\mathcal K)\) defined by \(\tilde u\) comes from the unit in weight zero, so it lands in \(F^0\), and its image in \(\operatorname{gr}^0\) is the tuple \((u_a)_a\). Composing with the filtered trace morphism of Lemma \hyperref[bp45]{11.16} and forgetting the filtration gives on the one hand \(\operatorname{Tr}_c(u)\), and on the other hand, by the factorization through \(\operatorname{gr}^0\), the sum of the trace classes \(\operatorname{Tr}_c(u_a)\). For (2) we apply (1) to the lift \(\tilde u\) of Lemma \hyperref[bp48]{11.17}.
\end{proof}

% text: prose-11-21
\phantomsection\label{bp56}%
Restricting the formula to a proper open and closed part \(B\) of the fixed locus and integrating, we get \(\operatorname{LT}_B(u)=\operatorname{LT}_B(u_U)+\operatorname{LT}_B(u_Z)\).



\clearpage\phantomsection\hypertarget{bp-outside}{}%
\section*{Outside this export}
\addcontentsline{toc}{section}{Outside this export}
References marked \ensuremath{{}^\ast} are to parts of the blueprint this document leaves out:
\begin{description}
\item[Section 3.2] Constructibility, in Ⅰ The six operations (\bpslug{sec-{}\allowbreak{}Constructibility})
\item[Lemma 3.13] Flat models and generators, in Ⅰ The six operations (\bpslug{lemma-{}\allowbreak{}flat-{}\allowbreak{}models-{}\allowbreak{}generators})
\item[Theorem 5.2] The exceptional inverse image, in Ⅰ The six operations (\bpslug{theorem-{}\allowbreak{}exceptional-{}\allowbreak{}inverse-{}\allowbreak{}image})
\item[Proposition 5.15] Comparison with the étale derived category, in Ⅰ The six operations (\bpslug{proposition-{}\allowbreak{}comparison-{}\allowbreak{}etale-{}\allowbreak{}derived-{}\allowbreak{}category})
\item[Proposition 7.1] The six operations on $\Dctf$, in Ⅰ The six operations (\bpslug{proposition-{}\allowbreak{}six-{}\allowbreak{}operations-{}\allowbreak{}dctf})
\item[Section 8.9] The Künneth formula, in Ⅰ The six operations (\bpslug{sec-{}\allowbreak{}The-{}\allowbreak{}Kunneth-{}\allowbreak{}formula})
\item[Proposition 8.18] Ordinary external Künneth, in Ⅰ The six operations (\bpslug{proposition-{}\allowbreak{}ordinary-{}\allowbreak{}external-{}\allowbreak{}kunneth})
\item[Theorem 9.3] Smooth purity for all admissible coefficients, in Ⅱ Duality and the trace formula (\bpslug{theorem-{}\allowbreak{}smooth-{}\allowbreak{}purity-{}\allowbreak{}all-{}\allowbreak{}admissible-{}\allowbreak{}coefficients})
\item[Definition 9.24] Trace map, in Ⅱ Duality and the trace formula (\bpslug{definition-{}\allowbreak{}trace-{}\allowbreak{}map})
\item[Definition 10.2] Verdier dual and evaluation, in Ⅱ Duality and the trace formula (\bpslug{definition-{}\allowbreak{}verdier-{}\allowbreak{}dual-{}\allowbreak{}evaluation})
\item[Corollary 10.9] Verdier exchange formulas, in Ⅱ Duality and the trace formula (\bpslug{corollary-{}\allowbreak{}verdier-{}\allowbreak{}exchange-{}\allowbreak{}formulas})
\item[Section 10.5] Weil complexes, in Ⅱ Duality and the trace formula (\bpslug{sec-{}\allowbreak{}Weil-{}\allowbreak{}complexes})
\item[Chapter 12] The trace formula, in Ⅱ Duality and the trace formula (\bpslug{chap-{}\allowbreak{}The-{}\allowbreak{}trace-{}\allowbreak{}formula})
\item[Definition 12.2] Frobenius correspondences and their operators, in Ⅱ Duality and the trace formula (\bpslug{definition-{}\allowbreak{}frobenius-{}\allowbreak{}correspondences-{}\allowbreak{}their-{}\allowbreak{}operators})
\item[Remark 12.4] Local terms are not computed stalk by stalk, in Ⅱ Duality and the trace formula (\bpslug{remark-{}\allowbreak{}locality-{}\allowbreak{}warning})
\item[Proposition 12.16] The Frobenius local term on a curve, in Ⅱ Duality and the trace formula (\bpslug{proposition-{}\allowbreak{}frobenius-{}\allowbreak{}local-{}\allowbreak{}term-{}\allowbreak{}curve})
\item[Section 12.5] $\ell$-adic coefficients, in Ⅱ Duality and the trace formula (\bpslug{sec-{}\allowbreak{}ell-{}\allowbreak{}adic-{}\allowbreak{}coefficients-{}\allowbreak{}5})
\end{description}
\clearpage\phantomsection%
\section*{References}
\addcontentsline{toc}{section}{References}
\bpbibitem{BBD82}{bp58}{A. Beilinson, J. Bernstein, and P. Deligne. \emph{Faisceaux pervers}, volume 100 of \emph{Astérisque}. 1982.}

\bpbibitem{BS15}{bp59}{B. Bhatt and P. Scholze. The pro-étale topology for schemes. \emph{Astérisque}, 369:99–201, 2015. \href{https://people.mpim-bonn.mpg.de/scholze/proetale.pdf}{Author's manuscript}.}

\bpbibitem{Con07}{bp60}{B. Conrad. Deligne's notes on Nagata compactifications. \emph{Journal of the Ramanujan Mathematical Society}, 22:205–257, 2007. \href{https://math.stanford.edu/\~conrad/papers/nagatafinal.pdf}{Author's manuscript}.}

\bpbibitem{Del74}{bp61}{P. Deligne. La conjecture de Weil. I. \emph{Publ. Math. IHÉS}, 43:273–307, 1974.}

\bpbibitem{Del80}{bp62}{P. Deligne. La conjecture de Weil. II. \emph{Publ. Math. IHÉS}, 52:137–252, 1980.}

\bpbibitem{Den22}{bp63}{C. Deninger. There is no “Weil-”cohomology theory with real coefficients for arithmetic curves. \emph{Ann. Sc. Norm. Super. Pisa Cl. Sci.}, 2022. To appear. arXiv:\href{https://arxiv.org/abs/2204.02714}{2204.02714}.}

\bpbibitem{Deu41}{bp64}{M. Deuring. Die Typen der Multiplikatorenringe elliptischer Funktionenkörper. \emph{Abh. Math. Sem. Univ. Hamburg}, 14:197–272, 1941.}

\bpbibitem{Eke90}{bp65}{T. Ekedahl. On the adic formalism. In \emph{The Grothendieck Festschrift, Vol. II}, volume 87 of \emph{Progress in Mathematics}, pages 197–218. Birkhäuser, 1990.}

\bpbibitem{Has36}{bp66}{H. Hasse. Zur Theorie der abstrakten elliptischen Funktionenkörper. I, II, III. \emph{J. Reine Angew. Math.}, 175:55–62, 69–88, 193–208, 1936.}

\bpbibitem{Hum72}{bp67}{J. E. Humphreys. \emph{Introduction to Lie Algebras and Representation Theory}, volume 9 of \emph{Graduate Texts in Mathematics}. Springer, 1972.}

\bpbibitem{Ill94}{bp68}{L. Illusie. Autour du théorème de monodromie locale. \emph{Astérisque}, 223:9–57, 1994. §4. \href{https://www.numdam.org/item/AST_1994__223__9_0/}{Original article}.}

\bpbibitem{Kat88}{bp69}{N. M. Katz. Travaux de Laumon. \emph{Astérisque}, 161–162:105–132, 1988. Séminaire Bourbaki, exp. 691.}

\bpbibitem{KW01}{bp70}{R. Kiehl and R. Weissauer. \emph{Weil conjectures, perverse sheaves and \(\ell\)-adic Fourier transform}, volume 42 of \emph{Ergebnisse der Mathematik}. Springer, 2001.}

\bpbibitem{Lau87}{bp71}{G. Laumon. Transformation de Fourier, constantes d'équations fonctionnelles et conjecture de Weil. \emph{Publ. Math. IHÉS}, 65:131–210, 1987.}

\bpbibitem{Mil86a}{bp72}{J. S. Milne. \href{https://www.jmilne.org/math/xnotes/AVs.pdf}{Abelian Varieties}. In \emph{Arithmetic Geometry}. 1986. Corrected online version.}

\bpbibitem{Mil86b}{bp73}{J. S. Milne. \href{https://www.jmilne.org/math/xnotes/JVs.pdf}{Jacobian Varieties}. In \emph{Arithmetic Geometry}. 1986. Corrected online version.}

\bpbibitem{Ran39}{bp74}{R. A. Rankin. Contributions to the theory of Ramanujan's function \(\tau(n)\) and similar arithmetical functions. II. \emph{Proc. Cambridge Philos. Soc.}, 35:357–372, 1939.}

\bpbibitem{SGA4}{bp75}{M. Artin, A. Grothendieck, and J.-L. Verdier, editors. \emph{Théorie des topos et cohomologie étale des schémas}, volumes 269, 270, 305 of \emph{Lecture Notes in Mathematics}. Springer, 1972–1973. SGA 4.}

\bpbibitem{SGA4½}{bp76}{P. Deligne. \emph{Cohomologie étale}, volume 569 of \emph{Lecture Notes in Mathematics}. Springer, 1977. SGA 4½; \emph{Théorèmes de finitude}, §3, and \emph{Rapport sur la formule des traces}, §3, and \emph{Sommes trigonométriques}, §1.}

\bpbibitem{SGA5}{bp57}{A. Grothendieck, editor. \emph{Cohomologie \(\ell\)-adique et fonctions \(L\)}, volume 589 of \emph{Lecture Notes in Mathematics}. Springer, 1977. SGA 5.}

\bpbibitem{SGA7II}{bp77}{P. Deligne and N. Katz, editors. \emph{Groupes de monodromie en géométrie algébrique}, volume 340 of \emph{Lecture Notes in Mathematics}. Springer, 1973. SGA 7 II; Exposé XIII, (2.1.7).}

\bpbibitem{Stacks}{bp78}{The Stacks Project Authors. \href{https://stacks.math.columbia.edu}{The Stacks Project}. Individual results are cited by their stable tags in the text.}

\bpbibitem{Var05}{bp79}{Y. Varshavsky. Lefschetz–Verdier trace formula and a generalization of a theorem of Fujiwara. 2005. Version 2. arXiv:\href{https://arxiv.org/abs/math/0505564}{math/0505564}. \href{https://arxiv.org/abs/math/0505564v2}{Author's preprint}.}

\bpbibitem{Wei48a}{bp80}{A. Weil. \emph{Sur les courbes algébriques et les variétés qui s'en déduisent}. Hermann, Paris, 1948.}

\bpbibitem{Wei48b}{bp81}{A. Weil. \emph{Variétés abéliennes et courbes algébriques}. Hermann, Paris, 1948.}

\bpbibitem{Wei49}{bp82}{A. Weil. Numbers of solutions of equations in finite fields. \emph{Bull. Amer. Math. Soc.}, 55:497–508, 1949.}

\end{document}
